%% 第24回編集委員会（物理第3巻）：組 波W3
%% p3-05.tex  波-5　波を式で書く：進行波と波動方程式（★2〜4・3段にしない）
%% 中身は計画書 tools/phys3-plan.md の B「波-5」。公式の持ち主は A（A4・A11〜A14・A19 の速さ・A26・A54）。
%% 第2部（波）の入口：第1部の「おもりとばね」を、自由度を無限にふやすと波になる（⑦。計画書 0 の3）。
%% 分散関係と群速度は、この章の⑦で作る（数学の本にない道具。C-6-1）。波-7 ⑥は、ここの群速度を指して水の波へ使う。
%% 数値はすべて Python で計算した：checks/物理第3巻/波-5/calc_p3_05.py
%% 記号（この章）：波の速さ c（W-3。第1部の減衰係数 c はこの章に出ない）、波長 λ、波数 k〔rad/m〕、角振動数 ω、振動数 f、周期 T、
%%   振幅 A（複素振幅だけ X）、位相 φ、形の関数 y₀（8-13 の f は振動数 f とぶつかるので使わない）、張力 F_T（8-13 は F）、線密度 μ、
%%   ヤング率 E、体積弾性率 K、密度 ρ、棒・管の中の変位 u・ξ、水面の高さ η、水の深さ h、ばね定数 k_s（⑦で波数 k と同じ章に出るため）、
%%   おもりの間かく a、位相速度 c_p、群速度 c_g、差分の比 s＝cΔt/Δx（8-13 と同じ）。熱-4 ⑥は音の速さを a と書いた（この章では c）。
%% 第1巻との約束：力-16 ⑦（c＝√(E/ρ)≈5100 m/s、いちばん速いゆれ ω＝2c÷要素の大きさ、CFL）→ ⑤⑦⑧・例題11。
%%   力-8 の道しるべ（波-5）→ ⑤。第1巻の巻末「波長（波-5）」「波数（波-5）」→ ②③。
%% 数学の本との約束：8-13⑧⑨⑩（波動方程式・ダランベール・波数）を物理の側から。8-13:559 の「物理シリーズ『波』」は『波動』の誤り（報告に書いた）。
%%   9-1⑨ の A弦 110 Hz・0.65 m・143 m/s（例題7）。4-9⑦ の積と和の式（⑦の群速度）。
%% 出典（本文で使った実機の値・事実。第24回の追記により、書くときに確かめた）
%%   ・スタジアムのウェーブ（約 12 m/s・1秒に約 20 席・幅 6〜12 m、ふつう時計回り、大きなスタジアムでは 25〜35 人で始まる、
%%       右ききとの関係は著者の推測）：I. Farkas, D. Helbing, T. Vicsek, Nature 419, 131–132 (2002)
%%       （https://nature.berkeley.edu/garbelottoat/wp-content/uploads/Farkas_et_al-2002-Nature.pdf、計画書 D〔確〕）。
%%       arXiv 版（https://arxiv.org/pdf/cond-mat/0210073）の本文で「14 のウェーブ、12 m/s（20 seats/s）、幅 6〜12 m、more frequently clockwise、
%%       not more than a few dozens of people、右ききの人が多いことによる知覚のかたより（推測）」を読んだ。25〜35 人は NPR・Mental Floss の記事で照合。
%%   ・鋼 E＝205 GPa・ρ＝7850 kg/m³・G＝79 GPa、アルミニウム合金 69 GPa・2700 kg/m³、空気 1.20 kg/m³（20 ℃）：第1巻の巻末の表
%%   ・水の体積弾性率 2.2 GPa（力-8・力-12）。水の音速：IAPWS-95（Python の iapws で計算）で 20 ℃ 1482 m/s・密度 998.2 kg/m³、8 ℃ 1439 m/s
%%   ・ギターの弦の張力：D'Addario EJ16（Light 12-53）の表で A弦（.042）28.93 lb＝129 N、1弦（.012）23.36 lb＝104 N、6本の合計 160.7 lb
%%       （販売店2か所の表で照合：yamahamusiclondon.com・fiddlershop.com）。線密度 6.3 g/m・0.57 g/m は、この張力と 0.65 m・110 Hz・330 Hz から逆算した値
%%   ・地震波の速さ P 6.0 km/s・S 3.5 km/s は教科書の近似。大森の式の比例定数は 6〜8 km/s（岩手県教育委員会の地学基礎の教材、
%%       大阪教育大附属天王寺の地学の教材で照合）
%%   ・S 波は液体の外核を通らない（USGS の shadow zone の解説 https://www.usgs.gov/media/images/shadowzonegif と、
%%       IRIS（米国の地震学の大学連合）の教材 https://www.iris.edu/hq/inclass/animation/206 で照合。usgs.gov は作業環境から届かず、検索の要旨で読んだ）
%%   ・1960年のチリ地震の津波が日本に届くまで：内閣府 災害教訓の継承に関する専門調査会 報告書（1960 チリ地震津波）第2章で「23時間後」、
%%       日本赤十字社愛知県支部・名古屋地方気象台の資料で「約1日」。本文は「約1日」
%%   ・コラドンとシュトゥルム（1826年10月から、ジュネーブ湖の Rolle と Thonon の間 約 14000 m、夜に火薬の光を合図、9 s より長く 9.5 s より短い）：
%%       Colladon の手記の英訳（arXiv:1710.03777 https://arxiv.org/pdf/1710.03777）。1435 m/s・水温 8 ℃ は DOSITS（dosits.org）と ob-ultrasound.net で照合
%%   ・ラッセルの孤立波（1834年8月、ユニオン運河、高さ約 40 cm、馬で 1〜2 マイル追った）：Heriot-Watt 大学（macs.hw.ac.uk）の解説と
%%       Union Canal Unlocked の解説で照合
%%   ・グリーンの法則（浅い所へ来た波の高さ ∝ h^(−1/4)）：線形の浅水波の理論の結果
%%   ・ニュートンの音速（1687年、約15 % 小さい）とラプラス（1816年）：熱-4 のコラムに合わせた
%% 「この章で導く中学・高校の公式」の校種は計画書 A の案。弦の速さの式（高校物理（発展））は〔要確認：学習指導要領・教科書で照合できなかった〕。
\makeatletter
% 人（ウェーブ）  \rxfperson{x}{立ち上がりの量 0〜1}{色}
\newcommand\rxfperson[3]{%
  \pgfmathsetmacro\rxf@h{9+6*#2}%
  \fill[bkGray!35] (#1-2.2,0) rectangle (#1+2.2,1.2);
  \draw[#3!80!black, line width=1.6pt, line cap=round] (#1,1.8) -- (#1,\rxf@h-1.4);
  \fill[#3!70] (#1,\rxf@h) circle (1.5);
  \ifdim #2pt>0.55pt \draw[#3!80!black, line width=1pt, line cap=round] (#1,\rxf@h-3.2) -- (#1-1.8,\rxf@h+1.6) (#1,\rxf@h-3.2) -- (#1+1.8,\rxf@h+1.6);\fi}
% 波の太い矢印  \rxfarrow{(始め)}{(終わり)}
\newcommand\rxfarrow[2]{\draw[bkRed!80!black, line width=1.4pt, -{Stealth[length=3mm, width=2.4mm]}] #1 -- #2;}
% おもり（ばねとおもりの列）  \rxfmass{(中心)}
\newcommand\rxfmass[1]{\filldraw[fill=bkBlue!35, draw=bkBlue!70!black, line width=0.6pt] ($#1+(-2.4,-2.4)$) rectangle ($#1+(2.4,2.4)$);}
% ばね  \rxfspring{(左)}{(右)}
\newcommand\rxfspring[2]{\draw[bkInk!80, line width=0.6pt, decorate, decoration={zigzag, segment length=1.6mm, amplitude=1.3mm, pre length=1mm, post length=1mm}] #1 -- #2;}
\makeatother

\phchapter{5}{波を式で書く：\\進行波と波動方程式}{スタジアムのウェーブでは、何が進んでいる？}{\bklevelrange{2}{4}}{波を式で書く：進行波と波動方程式}
\bkmichishirube{4-9,6-10,8-13,力-8,力-16,波-1,波-3}{波-6,波-7,波-8,波-9}

\begin{goal}
\begin{itemize}[label={\color{bkTheme}●}, leftmargin=1.6\zw, itemsep=0pt, topsep=1pt]
\item \bklevel{2}\enspace 波で進むのは「ゆれの形」だと、横波・縦波で説明できる。
\item \bklevel{2}\enspace 振幅・波長・周期・振動数と $c=f\lambda$ を使いこなせる。
\item \bklevel{3}\enspace 正弦波の式 $y=A\sin(\omega t-kx)$ からゆれを求められる。
\item \bklevel{4}\enspace 媒質の運動方程式から波動方程式と波の速さを導ける。
\item \bklevel{4}\enspace おもりの列の分散関係・群速度と、計算の刻みの条件がわかる。
\end{itemize}
\end{goal}

\phkoushiki{%
\phk{音は物の振動が伝わる。真空では伝わらない}{中学理科}{ph:波-5:baishitsu}
\phk{横波と縦波、振幅・波長・周期・振動数、$v=f\lambda$}{高校物理基礎}{ph:波-5:flambda}
\phk{重ね合わせの原理・波の独立性}{高校物理基礎}{ph:波-5:kasane}
\phk{正弦波の式 $y=A\sin2\pi\Bigl(\dfrac{t}{T}-\dfrac{x}{\lambda}\Bigr)$}{高校物理}{ph:波-5:seigen}
\phk{弦を伝わる波の速さ $c=\sqrt{F_{\mathrm T}/\mu}$}{高校物理（発展）}{ph:波-5:gen}
}

\begin{hatena}
\noindent\begin{minipage}[c]{0.46\linewidth}\raggedright
野球やサッカーのスタジアムで、観客が順に立ち上がってはすわる「ウェーブ」が、客席をぐるっと回っていくことがあります。遠くから見ると、大きな波が客席を走っていくように見えます。

でも、客席を走って移動している人は、1人もいません。では、ウェーブでは、いったい何が進んでいるのでしょう。
\end{minipage}\hfill
\begin{minipage}[c]{0.52\linewidth}\centering
\begin{tikzpicture}[x=1mm, y=1mm]
  \foreach \i in {0,...,11} {%
    \pgfmathsetmacro\rxflift{1/(1+((\i-4)^2)/1.5)^2}
    {\rxfperson{\i*5.4}{\rxflift}{bkBlue}}}
  \rxfarrow{(34,21)}{(52,21)}
  \node[font=\sffamily\footnotesize, anchor=south] at (43,22.5) {ウェーブの進む向き};
  \node[font=\sffamily\footnotesize, anchor=north] at (30,-1) {人は立ってすわるだけ};
\end{tikzpicture}
\end{minipage}
\end{hatena}

\noindent 第1部（波-1〜4）では、おもりとばねのゆれ（振動）を調べてきました。波-3 では、おもりがいくつもつながると、ゆれがとなりへ、またとなりへと伝わることを見ました。この章からの第2部では、おもりの数を限りなく増やして、ゆれが\textbf{伝わっていく}ようすそのもの、つまり\textbf{波}を調べます。弦も、鉄の棒も、水も、空気も、小さなおもりが限りなく並んで、となりどうしがばねでつながったものと考えられます。だから、ゆれの式から、波の式が作れるのです。

\begin{yattemiyou}[ばねのおもちゃで、波を作ろう]
\textbf{用意するもの}：ばねのおもちゃ（くるくるのびる大きなばね）、リボン。2人で行う\par\vspace{2pt}
\noindent\begin{minipage}[c]{0.56\linewidth}
\begin{enumerate}[label={\bkcircled[bkGreen]{\arabic*}}, leftmargin=2.4\zw, itemsep=1pt, topsep=0pt]
\item ばねを床に\qty{3}{m}くらいのばして両はしを持ち、まん中の輪にリボンを結ぶ。
\item 片方の人が、手を1回だけ横にふる。ゆれとリボンの動きを見る。ばねを強く引いてから、もう一度ふる。
\item 手を1回だけ前に押して、もどす。輪のつまった所がどう進むかを見る。
\end{enumerate}
\end{minipage}\hfill
\begin{minipage}[c]{0.42\linewidth}\centering
\begin{tikzpicture}[x=1mm, y=1mm]
  \draw[bkInk!70, line width=0.5pt] plot[smooth, domain=0:52, samples=120] (\x, {6/(1+((\x-30)/4)^2)});
  \foreach \xx in {0,1.6,...,52} {\pgfmathsetmacro\yy{6/(1+((\xx-30)/4)^2)}\draw[bkInk!75, line width=0.45pt] (\xx,\yy-1.6) -- (\xx,\yy+1.6);}
  \fill[bkRed] (17.6,0.4) circle (0.9);
  \node[font=\sffamily\footnotesize, anchor=north] at (17.6,-2) {リボン};
  \draw[bkRed!80!black, line width=1pt, -{Stealth[length=2mm]}] (34,9.5) -- (44,9.5);
  \node[font=\sffamily\footnotesize, anchor=south] at (34,10.5) {ゆれが進む};
  \begin{scope}[shift={(0,-16)}]
  \foreach \xx in {0,1.8,...,14,15.2,16.2,...,22,24,26,...,52} {\draw[bkInk!75, line width=0.45pt] (\xx,-1.6) -- (\xx,1.6);}
  \draw[bkRed!80!black, line width=1pt, -{Stealth[length=2mm]}] (20,4) -- (30,4);
  \node[font=\sffamily\footnotesize, anchor=north] at (19,-2.4) {つまった所が進む};
  \end{scope}
\end{tikzpicture}
\end{minipage}
\end{yattemiyou}

\noindent やってみようの答え：横にふると、こぶの形がばねにそって向こうへ走っていきます。リボンは、こぶが通るときに横へ動いて、またもとの場所にもどるだけで、こぶといっしょに向こうへ行ったりはしません。前に押すと、輪がつまった所が向こうへ走ります。このときリボンは、前後に少し動いてもどるだけです。ばねを強く引くと、ゆれは速く進みます。\textbf{進んでいくのは、ばねそのものではなく、ゆれの形}なのです。どうして強く引くと速くなるのかは、④で式から確かめます。

\begin{omoidasou}
\bkrecall{4-8}{角 $\theta=\omega t$。周期 $T=2\pi/\omega$、振動数 $f=1/T$〔Hz〕。}
\bkrecall{5-15}{$\sin\omega t$ を $t$ で2回微分すると、もとの $-\omega^2$ 倍。}
\bkrecall{8-13}{波動方程式 $\dfrac{\partial^2y}{\partial t^2}=c^2\dfrac{\partial^2y}{\partial x^2}$、$c=\sqrt{F/\mu}$（張力を $F$ と書いた）。$y=f(x-ct)$ は形を保って走る（この章では $y_0$ と書く）。波数 $k=2\pi/\lambda$、$\omega=ck$。}
\bkrecall{6-10}{$e^{j\theta}=\cos\theta+j\sin\theta$。$e^{j\omega t}$ を微分すると $j\omega$ がかかる。}
\bkrecall{力-8}{応力 $\sigma=E\varepsilon$（$E$ はヤング率）。体積弾性率 $K$。棒のばね定数 $EA/L$（この章では断面積を $S$ と書く）。}
\bkrecall{波-3}{おもりがいくつもつながると、ゆれは、となりへ順に伝わる。}
\end{omoidasou}

\noindent この章の記号：高校の教科書では、波の速さを $v$、音の速さを $V$ と書きます。この本では、波の速さを $c$ と書きます（数学 8-13・9-1、力-16 と同じ）。熱-4 ⑥は音の速さを $a$ と書きましたが、同じものです。また、数学 8-13 は弦の張力を $F$ と書きましたが、この本では、力の中でも「張力」であることがわかるように $F_{\mathrm T}$ と書きます。

%% ======================== 物理の道具 ========================
\bksection{doguu}{物理の道具}

\bksubsection{1}{波とは：形が進む}
\label{ph:波-5:sec1}
\bkhatenaref に答えましょう。ウェーブでは、1人1人は、その場で立ち上がって、すわるだけです。となりの人が立ち上がったのを見たら、少しおくれて自分も立ち上がる。これを順にくり返すので、「立ち上がっている」という\textbf{状態}が、客席を横へ横へと移っていきます。\textbf{ウェーブで進んでいるのは、人ではなく、「立ち上がっている」という形（状態）}です。これがはてな？の答えです。

調べた人たちの報告では、ウェーブは1秒に約20席、速さにして約\qty{12}{m/s}で進み、立ち上がっている部分の幅は\qtyrange{6}{12}{m}くらいです（\bkreidairef{1}）。100 m を約8秒で走る速さです。客席にすわったまま、陸上の選手と同じくらい速い「形」を作っているのです。

\begin{bkfigure}
\begin{tikzpicture}[x=1mm, y=1mm]
  \foreach \r/\sh/\lab in {0/0/{時刻 0}, 1/5/{0.25秒後}} {%
    \begin{scope}[shift={(0,-\r*25)}]
    \foreach \i in {0,...,16} {%
      \pgfmathsetmacro\rxflift{1/(1+((\i-4-\sh)^2)/1.5)^2}
      \ifnum\i=6 {\rxfperson{\i*5.4}{\rxflift}{bkRed}}\else {\rxfperson{\i*5.4}{\rxflift}{bkBlue}}\fi}
    \node[font=\sffamily\footnotesize, anchor=south west] at (-2,17.5) {\lab};
    \end{scope}}
  \draw[bkGray, line width=0.5pt, dashed] (21.6,17) -- (21.6,-27);
  \draw[bkGray, line width=0.5pt, dashed] (48.6,17) -- (48.6,-27);
  \draw[bkInk, line width=0.6pt, -{Stealth[length=2mm]}] (21.6,-28.5) -- (48.6,-28.5);
  \node[font=\sffamily\footnotesize, anchor=north] at (35.1,-29) {5席分（\qty{3}{m}）進む};
  \node[font=\sffamily\footnotesize, anchor=west, text width=30mm, align=left] at (94,-4) {左から7番めの人（色を変えた人）は、同じ席で立ったりすわったりするだけ};
\end{tikzpicture}
\bkfigcaption{ウェーブの形が進むようす。0.25秒のあいだに、立ち上がっている所が5席分（席の幅 \qty{0.6}{m} なら \qty{3}{m}）右へ移る。左から7番めの人（色を変えた人）は、席を動かない。}\label{ph:波-5:fig-wave}
\end{bkfigure}

\phboxhead{媒質と波}
ばねのおもちゃや、なわ、水、空気のように、波を伝える物を\textbf{媒質}（ばいしつ）といいます。\bkidx{ばいしつ}{媒質}最初にゆれを作る所を\textbf{波源}（はげん）といいます。ある所のゆれが、媒質の中を次々に伝わっていく現象を\textbf{波}（波動）といいます。\bkidx{なみ}{波（波動）}媒質の1つ1つの部分は、もとの場所の周りでゆれるだけで、波といっしょに遠くへ運ばれることはありません。

\phboxhead{横波と縦波}
\label{ph:波-5:yokotate}やってみようの②のように、媒質のゆれる向きが、波の進む向きに\textbf{垂直}な波を\textbf{横波}（よこなみ）といいます。③のように、ゆれる向きが、波の進む向きと\textbf{同じ}（平行）な波を\textbf{縦波}（たてなみ）といいます。\bkidx{よこなみ}{横波}\bkidx{たてなみ}{縦波}縦波では、媒質がつまった所（\textbf{密}）とまばらな所（\textbf{疎}）が交互にでき、それが進んでいくので、\textbf{疎密波}（そみつは）ともいいます。\bkidx{そみつは}{疎密波}音は、空気の疎密波です。

\begin{bkfigure}
\begin{tikzpicture}[x=1mm, y=1mm]
  % 横波
  \node[font=\sffamily\footnotesize, anchor=west] at (0,30) {横波：ゆれ（上下）は、進む向き（右）に垂直};
  \foreach \i in {0,...,24} {\pgfmathsetmacro\yy{5*sin(\i*30)}\fill[bkBlue!70] (\i*4.4,18+\yy) circle (1.1);}
  \draw[bkGray, line width=0.4pt, dashed] (0,18) -- (106,18);
  \draw[bkInk, line width=0.7pt, {Stealth[length=1.6mm]}-{Stealth[length=1.6mm]}] (13.2,12) -- (13.2,24);
  \rxfarrow{(110,18)}{(122,18)}
  % 縦波
  \node[font=\sffamily\footnotesize, anchor=west] at (0,7) {縦波：ゆれ（左右）は、進む向き（右）と同じ};
  \foreach \i in {0,...,24} {\pgfmathsetmacro\xx{\i*4.4+3*sin(\i*30)}\fill[bkBlue!70] (\xx,0) circle (1.1);}
  \draw[bkInk, line width=0.7pt, {Stealth[length=1.6mm]}-{Stealth[length=1.6mm]}] (8,-4.5) -- (18,-4.5);
  \rxfarrow{(110,0)}{(122,0)}
  \node[font=\sffamily\footnotesize, anchor=north] at (26.4,-3) {密};
  \node[font=\sffamily\footnotesize, anchor=north] at (52.8,-3) {疎};
  \node[font=\sffamily\footnotesize, anchor=north] at (79.2,-3) {密};
\end{tikzpicture}
\bkfigcaption{横波と縦波。点は媒質の小さな部分。縦波では、つまった所（密）とまばらな所（疎）が交互にできる。}\label{ph:波-5:fig-yokotate}
\end{bkfigure}

\phboxhead{縦波を横波のように書く}
縦波は、点の並びのままでは、どこがどれだけずれたかが読みにくいので、グラフにして表します。それぞれの点の、もとの位置からのずれ（右へのずれを正とする）を、その点の位置の上に\textbf{たての長さとして}書くのです（図\ref{ph:波-5:fig-tatehyouji}）。これを\textbf{縦波の横波表示}といいます。\bkidx{たてなみのよこなみひょうじ}{縦波の横波表示}グラフが0を\textbf{右下がりに}横切る所では、左の点が右へ、右の点が左へずれているので、点が集まって密になります。右上がりに横切る所は疎です。

\begin{bkfigure}
\begin{tikzpicture}[x=1mm, y=1mm]
  \foreach \i in {0,...,24} {%
    \pgfmathsetmacro\dx{3*sin(\i*30)}
    \draw[bkGray!70, line width=0.35pt] (\i*4.4,-8) -- (\i*4.4,-5.5);
    \fill[bkBlue!70] (\i*4.4+\dx,-8) circle (1.1);
    \draw[bkOrange!80!black, line width=0.6pt, -{Stealth[length=1.3mm]}] (\i*4.4,3) -- ++(0,{2*\dx});}
  \draw[bkInk, line width=0.5pt, -{Stealth[length=2mm]}] (-2,3) -- (112,3) node[right, font=\footnotesize] {$x$};
  \draw[bkOrange!85!black, line width=1.1pt, domain=0:105.6, samples=120, smooth] plot (\x, {3+6*sin(\x/4.4*30)});
  \node[font=\sffamily\footnotesize, anchor=south] at (26.4,10.5) {右へのずれを上に};
  \node[font=\sffamily\footnotesize] at (26.4,-12.8) {密};
  \node[font=\sffamily\footnotesize] at (52.8,-12.8) {疎};
  \node[font=\sffamily\footnotesize] at (79.2,-12.8) {密};
  \node[font=\sffamily\footnotesize, anchor=west, text width=16mm, align=left] at (110,-8) {点のずれ};
\end{tikzpicture}
\bkfigcaption{縦波の横波表示。下の点のずれ（右向きを正）を、上のグラフのたての長さで表す。グラフが0を右下がりに横切る所が密、右上がりに横切る所が疎。}\label{ph:波-5:fig-tatehyouji}
\end{bkfigure}

\phboxhead{地震の波：P波とS波}
地震のときに地面の中を伝わる波には、縦波と横波の両方があります。速い縦波は、はじめに届くので\textbf{P波}（Primary＝最初の）、おそい横波は2番めに届くので\textbf{S波}（Secondary＝2番めの）とよばれます。\bkidx{ぴーは}{P 波・S 波}地面の浅い所では、P波は1秒に約\qty{6}{km}、S波は約\qty{3.5}{km}進みます。P波が来ると、カタカタと小さくゆれ（初期微動）、S波が来ると、ユサユサと大きくゆれます（主要動）。2つの波の届く時間の差から、震源までの距離がわかります（\bkreidairef{6}）。

\phboxhead{液体や気体は、横波を伝えない}
横波は、となりの部分を\textbf{ずらす}ように動かします。固体は、ずらされると、もとにもどろうとする力（ずれの弾性。力-8 ④の横弾性係数）がはたらくので、横波を伝えます。液体や気体は、ずらしても、もどろうとする力がはたらかない（だから流れる）ので、横波を伝えません。押し縮めればもどろうとする（体積弾性率）ので、縦波は伝えます。地球の中心の近くには、液体の鉄でできた層（外核）があり、S波はそこを通りぬけられません。このことから、外核が液体であることがわかりました。

\phboxhead{音は、真空では伝わらない}
\label{ph:波-5:baishitsu}音は、物のゆれが、まわりの空気を押したり引いたりして、空気の疎密波として伝わるものです。ガラスの入れ物の中でブザーを鳴らし、ポンプで空気をぬいていくと、音はだんだん小さくなり、ほとんど聞こえなくなります。音を伝える媒質（空気）がないからです。このあと④⑤で、波の速さは「媒質のもどろうとする力」と「媒質の重さ」から決まることを、式で確かめます。媒質がなければ、もどろうとする力も重さもないので、音の式そのものが作れないのです（光は媒質がなくても伝わる特別な波です。物理シリーズの『電磁気学』（電-15））。

\begin{phmijika}[身の回りの波]
田んぼや麦畑を風がふきぬけると、稲の穂が順にかたむいて起き上がり、波の形が畑を走っていきます。稲は、根もとで止まったまま、ゆれているだけです。高速道路の渋滞も、波のように伝わります。前の車がブレーキをふむと、後ろの車が順にブレーキをふむので、「止まっている所」が、車の進む向きとは\textbf{反対に}、後ろへ後ろへと伝わっていきます。車は前に進んでいるのに、渋滞の波は後ろへ進むのです。
\end{phmijika}

\bksubsection{2}{波の要素と $c=f\lambda$}
\label{ph:波-5:sec2}
\label{ph:波-5:youso}なわの片はしを、手で同じリズムでくり返し上下にゆらし続けると、なわには、山と谷がくり返し並んだ波ができ、それが向こうへ進んでいきます。ある瞬間の、なわの形を写真にとったものが図\ref{ph:波-5:fig-youso}の上のグラフです。

\begin{bkfigure}
\begin{tikzpicture}[x=1mm, y=1mm]
  % ある瞬間の形
  \draw[bkInk, line width=0.6pt, -{Stealth[length=2mm]}] (0,0) -- (98,0) node[right, font=\footnotesize] {$x$};
  \draw[bkInk, line width=0.6pt, -{Stealth[length=2mm]}] (0,-11) -- (0,13) node[above, font=\footnotesize] {$y$};
  \draw[bkBlue, line width=1.4pt, domain=0:92, samples=160, smooth] plot (\x, {8*sin(\x*9)});
  \draw[bkInk, line width=0.5pt, {Stealth[length=1.6mm]}-{Stealth[length=1.6mm]}] (10,10.5) -- (50,10.5);
  \node[font=\sffamily\footnotesize, anchor=south] at (30,11) {波長 $\lambda$};
  \draw[bkGray, line width=0.4pt, dashed] (10,8) -- (10,10.5) (50,8) -- (50,10.5);
  \draw[bkGray, line width=0.4pt, dashed] (90,8) -- (98,8);
  \draw[bkInk, line width=0.5pt, {Stealth[length=1.6mm]}-{Stealth[length=1.6mm]}] (96,0) -- (96,8);
  \node[font=\sffamily\footnotesize, anchor=west] at (97,5) {振幅 $A$};
  \node[font=\sffamily\footnotesize, anchor=south] at (50,-13) {};
  \node[font=\sffamily\footnotesize, anchor=north] at (30,-8.6) {谷};
  \node[font=\sffamily\footnotesize, anchor=south west] at (51,8.4) {山};
  \node[font=\sffamily\footnotesize, anchor=west] at (101,-4) {ある瞬間の形};
  \node[font=\sffamily\footnotesize, anchor=west] at (101,-9) {（$y$-$x$ グラフ）};
  % 1つの場所のゆれ
  \begin{scope}[shift={(0,-34)}]
  \draw[bkInk, line width=0.6pt, -{Stealth[length=2mm]}] (0,0) -- (98,0) node[right, font=\footnotesize] {$t$};
  \draw[bkInk, line width=0.6pt, -{Stealth[length=2mm]}] (0,-11) -- (0,13) node[above, font=\footnotesize] {$y$};
  \draw[bkOrange, line width=1.4pt, domain=0:92, samples=160, smooth] plot (\x, {8*sin(\x*9)});
  \draw[bkInk, line width=0.5pt, {Stealth[length=1.6mm]}-{Stealth[length=1.6mm]}] (10,10.5) -- (50,10.5);
  \node[font=\sffamily\footnotesize, anchor=south] at (30,11) {周期 $T$};
  \draw[bkGray, line width=0.4pt, dashed] (10,8) -- (10,10.5) (50,8) -- (50,10.5);
  \node[font=\sffamily\footnotesize, anchor=west] at (103,8) {1つの場所の};
  \node[font=\sffamily\footnotesize, anchor=west] at (103,3) {ゆれ（$y$-$t$ グラフ）};
  \end{scope}
\end{tikzpicture}
\bkfigcaption{上：ある瞬間の、なわの形（横の軸は場所 $x$）。となりの山までの長さが波長 $\lambda$。下：なわの1つの点のゆれ（横の軸は時刻 $t$）。同じ形に見えるが、横の軸が違う。1回ゆれる時間が周期 $T$。}\label{ph:波-5:fig-youso}
\end{bkfigure}

\begin{itemize}[label={\color{bkTheme}●}, leftmargin=1.6\zw, itemsep=1pt, topsep=2pt]
\item \textbf{振幅}（しんぷく）$A$〔m〕：ゆれていないときの位置から、山のいちばん高い所までの高さ（波-1 の単振動の振幅と同じ）。\bkidx{しんぷく}{振幅}
\item \textbf{波長}（はちょう）$\lambda$（ラムダ）〔m〕：となりどうしの山から山（谷から谷）までの長さ。\bkidx{はちょう}{波長}
\item \textbf{周期} $T$〔s〕：媒質の1つの点が、1回ゆれるのにかかる時間。
\item \textbf{振動数} $f$〔Hz〕：媒質の1つの点が、1秒にゆれる回数。$f=1/T$（波-1・4-8）。
\end{itemize}
上のグラフ（ある瞬間の形）の横の軸は\textbf{場所}、下のグラフ（1つの点のゆれ）の横の軸は\textbf{時刻}です。形は同じ $\sin$ に見えますが、山と山の間かくは、上では波長 $\lambda$、下では周期 $T$ で、意味がまったく違います。グラフを見たら、まず横の軸を確かめましょう。

\phboxhead{波の速さ：1周期で1波長進む}
手が1回ゆれる（時間 $T$）あいだに、なわには山が1つ、新しく作られて送り出されます。そのあいだに、1つ前に作った山は、ちょうど1波長 $\lambda$ だけ先へ進んでいます。時間 $T$ で $\lambda$ 進むので、波の速さは
\pheq{c=\frac{\lambda}{T}=f\lambda}{m/s ＝ Hz × m}
\label{ph:波-5:flambda}です（Hz は 1/s）。これが「波の速さ＝振動数×波長」です。高校の教科書の $v=f\lambda$ と同じ式です。

\begin{phyomi}[$c=f\lambda$]
「1秒に $f$ 個の山が送り出され、山と山の間かくが $\lambda$ なので、1秒に $f\lambda$ だけ波の先が進む」。$c$ は媒質で決まり（④⑤）、$f$ は波源（ゆらす手やスピーカー）で決まります。$\lambda$ は、その2つから決まる量です。
\end{phyomi}

\begin{phbaisuru}{振動数 $f$}
同じなわ（同じ媒質）なら、波の速さ $c$ は変わりません。手を2倍速くゆらすと、同じ時間に2倍の数の山が送り出されるので、山どうしがつまって、波長は\textbf{半分}になります。
\par\smallskip\noindent\begin{tikzpicture}[x=1mm, y=1mm]
  \draw[bkInk, line width=0.5pt, -{Stealth[length=1.6mm]}] (0,0) -- (104,0) node[right, font=\footnotesize] {$x$};
  \draw[bkBlue, line width=1.2pt, domain=0:96, samples=160, smooth] plot (\x, {5*sin(\x*11.25)});
  \node[font=\sffamily\footnotesize, anchor=west] at (98,5) {$f$};
  \draw[bkInk, line width=0.5pt, -{Stealth[length=1.6mm]}] (0,-16) -- (104,-16) node[right, font=\footnotesize] {$x$};
  \draw[bkOrange, line width=1.2pt, domain=0:96, samples=200, smooth] plot (\x, {-16+5*sin(\x*22.5)});
  \node[font=\sffamily\footnotesize, anchor=west] at (98,-11) {$2f$};
  \draw[bkInk, line width=0.4pt, {Stealth[length=1.4mm]}-{Stealth[length=1.4mm]}] (8,7.5) -- (40,7.5);
  \node[font=\sffamily\footnotesize, anchor=south, fill=white, inner sep=1pt] at (24,8) {$\lambda$};
  \draw[bkInk, line width=0.4pt, {Stealth[length=1.4mm]}-{Stealth[length=1.4mm]}] (4,-8.5) -- (20,-8.5);
  \node[font=\sffamily\footnotesize, anchor=south, fill=white, inner sep=1pt] at (12,-8) {$\lambda/2$};
\end{tikzpicture}
\end{phbaisuru}

\begin{phmijika}[音の高さと波長]
音の高さは振動数で決まり、振動数が大きいほど高い音です。人の耳に聞こえるのは、約\qty{20}{Hz}から約\qty{20}{kHz}までです。空気の中の音の速さ（\qty{20}{\celsius}で約\qty{343}{m/s}、熱-4 ⑥）で波長を計算すると、\qty{20}{Hz}の低い音は約\qty{17}{m}、\qty{20}{kHz}の高い音は約\qty{17}{mm}です（\bkreidairef{2}）。ピアノの低い音は、部屋の大きさほどもある波なのです。また、大きい音ほど、空気のゆれの振幅が大きくなっています（くわしくは波-9）。
\end{phmijika}

\bksubsection{3}{進む波の式}
\label{ph:波-5:sec3}
波の形を、式で書きましょう。時刻 $t=0$ の形が、関数 $y=y_0(x)$ で表されているとします。この形が、形を変えずに速さ $c$ で右へ進むと、時刻 $t$ には、グラフ全体が右へ $ct$ だけずれます。3-9 で、$x$ を $x-p$ にするとグラフは右へ $p$ だけずれることを学びました。だから、時刻 $t$ の形は
\pheq{y=y_0(x-ct)}{m}
です。左へ進む波なら $y=y_0(x+ct)$ です（8-13⑨のダランベールの解の2つの部分）。

\begin{bkfigure}
\begin{tikzpicture}
\begin{axis}[bkplot, width=120mm, height=40mm, xmin=-0.15, xmax=2.3, ymin=-0.02, ymax=0.14, xtick={0.5,1,1.5,2}, ytick={0.05,0.1}, yticklabel style={/pgf/number format/fixed}, xlabel={$x$（m）}, ylabel={$y$（m）}]
  \addplot[bkcurve, line width=1.3pt] coordinates {(-0.1,0) (0,0) (0.3,0.1) (0.4,0) (2.3,0)};
  \addplot[bkcurve2, line width=1.3pt, dash pattern=on 4pt off 2pt] coordinates {(0.5,0) (0.8,0.1) (0.9,0) (2.3,0)};
  \addplot[bkGreen!60!black, line width=1.3pt, dash pattern=on 1.2pt off 1.4pt] coordinates {(1.0,0) (1.3,0.1) (1.4,0)};
  \node[font=\sffamily\footnotesize, anchor=south] at (axis cs:0.3,0.103) {$t=0$};
  \node[font=\sffamily\footnotesize, anchor=south] at (axis cs:0.8,0.103) {0.25 s};
  \node[font=\sffamily\footnotesize, anchor=south] at (axis cs:1.3,0.103) {0.5 s};
  \draw[bkInk, line width=0.5pt, -{Stealth[length=1.6mm]}] (axis cs:1.5,0.06) -- (axis cs:2.0,0.06);
  \node[font=\sffamily\footnotesize, anchor=south] at (axis cs:1.75,0.064) {$c=\qty{2}{m/s}$};
\end{axis}
\end{tikzpicture}
\bkfigcaption{形を変えずに右へ進むパルス（1回だけのゆれ）。実線（$t=0$）、破線（\qty{0.25}{s}後）、点線（\qty{0.5}{s}後）。$c=\qty{2}{m/s}$ なら、0.25秒ごとに\qty{0.5}{m}ずつ右へずれる。}\label{ph:波-5:fig-pulse}
\end{bkfigure}

\phboxhead{正弦波の式}
\label{ph:波-5:seigen}波源（$x=0$）が、単振動 $y=A\sin\omega t$ でゆれ続けているとします（波-1）。その波が右へ速さ $c$ で進むと、場所 $x$ の点は、波源のゆれを、波が届くのにかかる時間 $x/c$ だけ\textbf{おくれて}くり返します。だから、場所 $x$・時刻 $t$ のずれは、時刻を $x/c$ だけもどした
\[y=A\sin\omega\Bigl(t-\frac{x}{c}\Bigr)=A\sin\Bigl(\omega t-\frac{\omega}{c}x\Bigr)\]
です。ここで、$\omega/c=2\pi f/(f\lambda)=2\pi/\lambda$ を $k$ と書くと、
\pheq{y=A\sin(\omega t-kx)\qquad k=\frac{2\pi}{\lambda},\quad\omega=\frac{2\pi}{T}}{m}
となります。この形の波を\textbf{正弦波}（せいげんは）といいます。\bkidx{せいげんは}{正弦波}$k$ は、8-13⑩の\textbf{波数}（はすう）で、1 m 進むごとに角が何ラジアン変わるかを表す数〔rad/m〕です。\bkidx{はすう}{波数}$\omega$ と $T$ の関係（$\omega T=2\pi$）と、$k$ と $\lambda$ の関係（$k\lambda=2\pi$）が、時間と場所で対になっています。$\omega=2\pi/T$、$k=2\pi/\lambda$ を入れると、
\[y=A\sin2\pi\Bigl(\frac{t}{T}-\frac{x}{\lambda}\Bigr)\]
とも書けます。高校の教科書の正弦波の式は、この形です。

\begin{phyomi}[$y=A\sin(\omega t-kx)$]
「場所 $x$ の点は、波源と同じ単振動をするが、その角が $kx$ だけおくれている」。$x$ が1波長 $\lambda$ 増えるごとに、おくれは $k\lambda=2\pi$（ちょうど1回分）なので、1波長はなれた点どうしは、まったく同じ動きをします。$\sin$ の中の $-kx$ のマイナスが「右へ進む」しるしで、左へ進む波は $A\sin(\omega t+kx)$ です。また、$\omega t-kx=-k(x-ct)$（$\omega=ck$）なので、正弦波も $y_0(x-ct)$ の形の1つです。
\end{phyomi}

\phboxhead{位相}
$\sin$ の中の角 $\omega t-kx$ を、\textbf{位相}（いそう）といいます（波-1 の単振動の位相と同じ）。\bkidx{いそう}{位相}位相が等しい2点は、同じ高さで同じ向きに動いていて、\textbf{同位相}（どういそう）といいます。位相が $\pi$（半周期分）違う2点は、いつも反対向きにずれていて、\textbf{逆位相}（ぎゃくいそう）といいます。1波長はなれた2点は同位相、半波長はなれた2点は逆位相です。ゆれ始めのずれを入れるときは、$y=A\sin(\omega t-kx+\varphi)$ と、初期位相 $\varphi$ をたします。

\phboxhead{媒質の速さは、波の速さと違う}
媒質の点が上下に動く速さは、$y$ を $t$ で微分した
\[\frac{\partial y}{\partial t}=A\omega\cos(\omega t-kx)\]
で（$x$ を止めて $t$ だけで微分するので、8-3 の偏微分の記号 $\partial$ を使う）、いちばん大きいときで $A\omega$ です。これは、波の形が進む速さ $c=\omega/k$ とは、まったく別の量です。振幅を2倍にすると、媒質の点は2倍速く動きますが、波が進む速さ $c$ は変わりません（\bkreidairef{3}）。

\begin{bkfigure}
\begin{tikzpicture}
\begin{axis}[bkplot, width=120mm, height=42mm, xmin=0, xmax=1.05, ymin=-0.03, ymax=0.032, xtick={0.2,0.4,0.6,0.8,1}, ytick={-0.02,0.02}, yticklabel style={/pgf/number format/fixed}, xticklabel style={/pgf/number format/fixed}, xlabel={$x$（m）}, ylabel={$y$（m）}, axis x line*=bottom]
  \draw[bkInk, line width=0.5pt] (axis cs:0,0) -- (axis cs:1.05,0);
  \addplot[bkcurve, domain=0:1.0, samples=160] {0.02*sin(deg(-5*pi*x))};
  \addplot[bkcurve2, dash pattern=on 4pt off 2pt, domain=0:1.0, samples=160] {0.02*sin(deg(10*pi*0.02-5*pi*x))};
  \fill[bkRed] (axis cs:0.2,0) circle (1.8pt);
  \draw[bkRed!80!black, line width=1pt, -{Stealth[length=1.6mm]}] (axis cs:0.2,0.002) -- (axis cs:0.2,0.017);
  \fill[bkRed] (axis cs:0.4,0) circle (1.8pt);
  \draw[bkRed!80!black, line width=1pt, -{Stealth[length=1.6mm]}] (axis cs:0.4,-0.002) -- (axis cs:0.4,-0.017);
  \node[font=\sffamily\footnotesize, anchor=south west, fill=white, inner sep=1pt] at (axis cs:0.62,0.021) {少しあとの形（破線）};
\end{axis}
\end{tikzpicture}
\bkfigcaption{右へ進む正弦波（実線）と、少しあとの形（破線。波は右へずれる）。媒質の点がどちらへ動くかは、少しあとの形を見ればわかる。$x=\qty{0.2}{m}$ の点は上へ、$x=\qty{0.4}{m}$ の点は下へ動いている。}\label{ph:波-5:fig-ugoki}
\end{bkfigure}

\bksubsection{4}{弦の運動方程式から波動方程式へ}
\label{ph:波-5:sec4}
\label{ph:波-5:gen}①〜③では、「波は形を保って速さ $c$ で進む」として式を作りました。では、その速さ $c$ は何で決まるのでしょう。ぴんと張った弦（ギターの弦、なわ）で、媒質の運動方程式（力-2）から確かめます。8-13⑧では、弦を小さな玉の列と考えて同じ式を作りました。ここでは、弦の短い一部分を、そのまま取り出して考えます。

\begin{bkfigure}
\begin{tikzpicture}[x=1mm, y=1mm]
  \draw[bkGray, line width=0.4pt, dashed] (0,0) -- (118,0);
  \draw[bkInk!50, line width=0.8pt] plot[smooth, domain=0:118, samples=80] (\x, {6+16*exp(-((\x-62)^2)/900)});
  \draw[bkBlue!80!black, line width=2pt] plot[smooth, domain=34:62, samples=40] (\x, {6+16*exp(-((\x-62)^2)/900)});
  \pgfmathsetmacro\yA{6+16*exp(-((34-62)^2)/900)}
  \fill[bkBlue!80!black] (34,\yA) circle (0.9) (62,22) circle (0.9);
  \draw[bkGray, line width=0.4pt] (34,\yA) -- ++(-16,0);
  \draw[bkRed!80!black, line width=1.2pt, -{Stealth[length=2.6mm]}] (34,\yA) -- ++(-157.4:14);
  \node[font=\sffamily\footnotesize, anchor=north east, text=bkRed!80!black] at (21,6.6) {$F_{\mathrm T}$};
  \draw[bkInk, line width=0.5pt] (24,\yA) arc (180:202.6:10);
  \node[font=\sffamily\footnotesize, anchor=south east] at (21,\yA) {$\theta_1$};
  \draw[bkRed!80!black, line width=1.2pt, -{Stealth[length=2.6mm]}] (62,22) -- ++(0:16);
  \node[font=\sffamily\footnotesize, anchor=south, text=bkRed!80!black] at (76,22.6) {$F_{\mathrm T}$};
  \draw[bkInk, line width=0.5pt] (34,-2) -- (34,-4) (62,-2) -- (62,-4);
  \draw[bkInk, line width=0.5pt, {Stealth[length=1.6mm]}-{Stealth[length=1.6mm]}] (34,-3) -- (62,-3);
  \node[font=\sffamily\footnotesize, anchor=north] at (34,-4.2) {$x$};
  \node[font=\sffamily\footnotesize, anchor=north] at (62,-4.2) {$x+\Delta x$};
  \node[font=\sffamily\footnotesize, anchor=north] at (48,-1) {$\Delta x$};
  \node[font=\sffamily\footnotesize, anchor=south east, text=bkBlue!80!black, align=right] at (47,21) {質量 $\mu\,\Delta x$};
\end{tikzpicture}
\bkfigcaption{弦の短い一部分（太い線、長さ $\Delta x$）にはたらく力。両はしを、となりの弦が張力 $F_{\mathrm T}$ で、弦にそった向きに引く。左はしの傾きの角は $\theta_1$、右はしの角 $\theta_2$ はここでは0。かたむきが違うので、たての力が残る。}\label{ph:波-5:fig-genbubun}
\end{bkfigure}

\noindent 弦の1 m あたりの質量を\textbf{線密度}（せんみつど）$\mu$（ミュー）〔kg/m〕とします。\bkidx{せんみつど}{線密度}$x$ から $x+\Delta x$ までの短い部分の質量は $\mu\,\Delta x$ です。この部分は、両はしを、となりの弦に張力 $F_{\mathrm T}$ で引かれます。左はしの弦のかたむきの角を $\theta_1$、右はしを $\theta_2$ とすると、たての力の合計は
\[F_{\mathrm T}\sin\theta_2-F_{\mathrm T}\sin\theta_1\]
です。かたむきが小さい（角が小さい）とき、$\sin\theta\approx\tan\theta$ で（5-17）、$\tan\theta$ はグラフの傾き $\partial y/\partial x$ です。だから
\begin{bkhenkei}[0.56]
\text{たての力}\approx F_{\mathrm T}\Bigl\{\frac{\partial y}{\partial x}(x+\Delta x)-\frac{\partial y}{\partial x}(x)\Bigr\} & $\sin\theta\approx\tan\theta$＝傾き \\
\approx F_{\mathrm T}\,\frac{\partial^2y}{\partial x^2}\,\Delta x & 傾きの差＝傾きの変わり方×$\Delta x$（5-3） \\
\mu\,\Delta x\,\frac{\partial^2y}{\partial t^2}=F_{\mathrm T}\,\frac{\partial^2y}{\partial x^2}\,\Delta x & 運動方程式（質量×加速度＝力） \\
\end{bkhenkei}
両辺を $\mu\,\Delta x$ でわると、
\pheq{\frac{\partial^2y}{\partial t^2}=c^2\,\frac{\partial^2y}{\partial x^2}\qquad c=\sqrt{\frac{F_{\mathrm T}}{\mu}}}{m/s ＝ $\sqrt{\text{N}\div\text{kg/m}}$}
となります。8-13⑧と同じ\textbf{波動方程式}（はどうほうていしき）です。\bkidx{はどうほうていしき}{波動方程式}横の向きの力は $F_{\mathrm T}\cos\theta_2-F_{\mathrm T}\cos\theta_1$ ですが、角が小さいと $\cos\theta\approx1$ なので、ほぼ0です。だから弦の各部分は、横へは動かず、たてにだけゆれます。

単位を確かめましょう。$\text{N}\div(\text{kg/m})=(\text{kg·m/s}^2)\times(\text{m/kg})=\text{m}^2/\text{s}^2$ なので、$\sqrt{F_{\mathrm T}/\mu}$ は確かに速さ〔m/s〕です。8-13⑨で確かめたように、どんな形 $y_0$ でも $y=y_0(x-ct)$ はこの式をみたすので、$c$ は、弦の上を形が進む速さそのものです。

\begin{phyomi}[$\mu\dfrac{\partial^2y}{\partial t^2}=F_{\mathrm T}\dfrac{\partial^2y}{\partial x^2}$]
「弦の曲がり方（$\partial^2y/\partial x^2$）が、その場所の加速度を決める」。下に凸に曲がった所（谷の底）は上向きに、上に凸に曲がった所（山の頂上）は下向きに加速されます。曲がりを元にもどそうとするのが張力、動きにくさが線密度です。だから\textbf{張力が大きく、軽い弦ほど、波は速く進みます}。
\end{phyomi}

\begin{phflow}[弦の波]
\phflowrow{$\mu\dfrac{\partial^2y}{\partial t^2}=F_{\mathrm T}\dfrac{\partial^2y}{\partial x^2}$}{かたむきが小さい\\張力は一定}{$c=\sqrt{\dfrac{F_{\mathrm T}}{\mu}}$}
\phflowrow{$\dfrac{\partial^2y}{\partial t^2}=c^2\dfrac{\partial^2y}{\partial x^2}$}{$y=A\sin(\omega t-kx)$\\を入れる}{$\omega=ck$、$c=f\lambda$}
\end{phflow}
2段めは、正弦波の式を波動方程式に入れると、$\partial^2y/\partial t^2=-\omega^2y$、$\partial^2y/\partial x^2=-k^2y$（5-15）なので、$-\omega^2y=c^2(-k^2y)$、つまり $\omega=ck$ になる、ということです。両辺を $2\pi$ でわると $f=c/\lambda$、すなわち②の $c=f\lambda$ です。②で「1周期で1波長」と考えて作った式が、運動方程式から出てきました。

\begin{phbaisuru}[4倍]{張力 $F_{\mathrm T}$}
$c=\sqrt{F_{\mathrm T}/\mu}$ なので、波の速さは $\sqrt4=2$ 倍です。張力を2倍にしただけでは、$\sqrt2\approx1.41$ 倍にしかなりません。線密度 $\mu$ を4倍（同じ材料で直径を2倍）にすると、速さは半分です。
\par\smallskip\noindent\begin{tikzpicture}
\begin{axis}[bkplot, width=100mm, height=42mm, xmin=0, xmax=560, ymin=0, ymax=330, xtick={129,258,516}, ytick={143,286}, xlabel={$F_{\mathrm T}$（N）}, ylabel={$c$（m/s）}]
  \addplot[bkcurve, domain=0:550, samples=100] {sqrt(x/0.0063)};
  \fill[bkRed] (axis cs:129,143.1) circle (1.8pt) (axis cs:516,286.2) circle (1.8pt);
  \draw[bkGray, line width=0.4pt, dashed] (axis cs:129,0) -- (axis cs:129,143.1) (axis cs:516,0) -- (axis cs:516,286.2);
  \node[font=\sffamily\footnotesize, anchor=north west, fill=white, inner sep=1pt] at (axis cs:280,150) {張力4倍で速さ2倍};
  \node[font=\sffamily\footnotesize, anchor=south east, text=bkGray, fill=white, inner sep=1pt] at (axis cs:550,20) {$\mu=\qty{6.3}{g/m}$ の弦（A弦）};
\end{axis}
\end{tikzpicture}
\end{phbaisuru}

\begin{phmodel}[弦の波動方程式]
\phignore{弦の重さ（重力）。空気の抵抗。弦を曲げるのに要る力（曲げのこわさ）。}
\phassume{弦のかたむきは小さい（$\sin\theta\approx\tan\theta$、$\cos\theta\approx1$）。ゆれても張力は変わらない。線密度はどこも同じ。}
\phlimit{大きくはじいたとき（弦がのびて張力が変わる）。太くてかたい弦（ピアノの低い音の弦は、曲げのこわさで倍音が少し高くなる。波-6）。}
\end{phmodel}

\begin{phmijika}[ギターのチューニング]
ギターの弦の音を高くするときは、糸巻き（ペグ）を回して弦を強く張ります。張力が大きくなると波が速くなり（$c$ が大きい）、同じ長さの弦の振動数が上がるからです（くわしくは波-6）。また、低い音の弦ほど太く、金属の線を巻いて重くしてあります。線密度 $\mu$ を大きくして、同じくらいの張力のまま、波をおそくしているのです。ふつうのアコースティックギターの弦では、6本の張力は、どれも\qtyrange{100}{130}{N}くらい（約\qtyrange{10}{13}{kg}のおもりの重さ）にそろえてあり、ネックが片側に曲がらないようになっています。
\end{phmijika}


\bksubsection{5}{棒・液体・気体・浅い水の波}
\label{ph:波-5:sec5}
弦では、もとにもどそうとする力が張力でした。棒・液体・気体・水面でも、同じように「小さな部分の運動方程式」を立てると、どれも同じ形の波動方程式になります。違うのは、もどそうとする力の強さと、重さの表し方だけです。

\phboxhead{棒を伝わる縦波}
細長い棒の一方のはしを、ハンマーでたたくと、「押された」という知らせが、縦波として棒を伝わります（力-16 ⑦の衝突の計算で約束した波です）。棒の各部分が、棒の向きにずれる量を $u(x,\ t)$ とします。

\begin{bkfigure}
\begin{tikzpicture}[x=1mm, y=1mm]
  \fill[bkGray!20] (0,0) rectangle (110,12);
  \draw[bkInk, line width=0.6pt] (0,0) rectangle (110,12);
  \fill[bkBlue!25] (40,0) rectangle (66,12);
  \draw[bkInk, line width=0.5pt, dashed] (36,-3) -- (36,15) (62,-3) -- (62,15);
  \draw[bkBlue!80!black, line width=1pt] (40,0) -- (40,12) (66,0) -- (66,12);
  \draw[bkInk, line width=0.5pt, -{Stealth[length=1.6mm]}] (36,13.5) -- (40,13.5);
  \node[font=\sffamily\footnotesize, anchor=south] at (37,14.2) {$u(x)$};
  \draw[bkInk, line width=0.5pt, -{Stealth[length=1.6mm]}] (62,13.5) -- (66,13.5);
  \node[font=\sffamily\footnotesize, anchor=south] at (67,14.2) {$u(x+\Delta x)$};
  \draw[bkRed!80!black, line width=1.2pt, -{Stealth[length=2.4mm]}] (40,6) -- (28,6);
  \node[font=\sffamily\footnotesize, anchor=south, text=bkRed!80!black] at (31,6.6) {$\sigma(x)S$};
  \draw[bkRed!80!black, line width=1.2pt, -{Stealth[length=2.4mm]}] (66,6) -- (80,6);
  \node[font=\sffamily\footnotesize, anchor=south, text=bkRed!80!black] at (78,6.6) {$\sigma(x+\Delta x)S$};
  \node[font=\sffamily\footnotesize, anchor=north] at (36,-3.3) {$x$};
  \node[font=\sffamily\footnotesize, anchor=north] at (62,-3.3) {$x+\Delta x$};
  \node[font=\sffamily\footnotesize, anchor=north] at (53,-7.6) {破線はずれる前、色の所はずれた後の小さな部分};
\end{tikzpicture}
\bkfigcaption{棒の小さな部分（断面積 $S$、長さ $\Delta x$。この章では振幅を $A$ と書くので、断面積は $S$ と書く）。左右の面を、となりの部分が応力 $\sigma$ で引く（押すときは $\sigma$ が負）。両はしのずれ $u$ が違うと、部分が伸び縮みする。}\label{ph:波-5:fig-bou}
\end{bkfigure}

\noindent 部分の伸びは $u(x+\Delta x)-u(x)$ なので、ひずみ（伸び÷もとの長さ、力-8）は $\partial u/\partial x$、応力は $\sigma=E\,\partial u/\partial x$ です。部分を右へ引く力と左へ引く力の差から、
\begin{bkhenkei}[0.56]
\rho S\,\Delta x\,\frac{\partial^2u}{\partial t^2}=\sigma(x+\Delta x)S-\sigma(x)S & 運動方程式（質量 $\rho S\,\Delta x$） \\
=ES\Bigl\{\frac{\partial u}{\partial x}(x+\Delta x)-\frac{\partial u}{\partial x}(x)\Bigr\} & $\sigma=E\,\partial u/\partial x$ \\
\approx ES\,\frac{\partial^2u}{\partial x^2}\,\Delta x & ④と同じ「傾きの差」 \\
\end{bkhenkei}
両辺を $\rho S\,\Delta x$ でわると、$\partial^2u/\partial t^2=(E/\rho)\,\partial^2u/\partial x^2$。弦と同じ形の波動方程式で、
\pheq{c=\sqrt{\frac{E}{\rho}}}{m/s ＝ $\sqrt{\text{Pa}\div\text{kg/m}^3}$}
です。鋼（$E=\qty{205}{GPa}$、$\rho=\qty{7850}{kg/m^3}$）では、
\[c=\sqrt{\frac{205\times10^9}{7850}}\approx\qty{5100}{m/s}\]
で、力-16 ⑦で使った値です。空気の中の音の約15倍の速さです。

\phboxhead{液体・気体を伝わる縦波}
管に入った水や空気でも、同じように考えます。小さな部分のずれを $\xi$（グザイ）とすると、部分の体積が変わる割合は $\partial\xi/\partial x$ で、圧力は、体積弾性率 $K$（力-8 ④）で $K\times$（縮んだ割合）だけ増えます。
\[\text{増えた圧力}\ p'=-K\,\frac{\partial\xi}{\partial x}\]
（縮むと $\partial\xi/\partial x<0$ で、圧力が増えるのでマイナス）。左の面は右へ $p'(x)S$、右の面は左へ $p'(x+\Delta x)S$ で押されるので、棒の $\sigma$ を $-p'$ にかえた式になり、
\pheq{c=\sqrt{\frac{K}{\rho}}}{m/s ＝ $\sqrt{\text{Pa}\div\text{kg/m}^3}$}
となります。水（$K=\qty{2.2}{GPa}$、$\rho=\qty{998}{kg/m^3}$）では約\qty{1500}{m/s}です（\bkreidairef{9}）。

気体では、$K$ に何を入れるかが大切です。ゆっくり押して温度が変わらないときは $K=p$（力-12）、すばやく押して熱が逃げないときは $K_{\mathrm s}=\gamma p$（熱-3 ⑧）でした。音のゆれは速いので $K_{\mathrm s}$ を使い、
\[c=\sqrt{\frac{\gamma p}{\rho}}\]
となります。熱-4 ⑥で、波の先をかこむ検査体積から導いた式（$a=\sqrt{\gamma p/\rho}$）と同じ結果が、運動方程式からも出てきました。\qty{20}{\celsius}の空気で約\qty{343}{m/s}です。気体の音の速さが温度でどう変わるかは、波-9 でくわしく学びます。

\begin{phyomi}[$c=\sqrt{\text{もどそうとする力の強さ}\div\text{重さ}}$]
弦・棒・液体・気体の波の速さは、どれも「$\sqrt{\text{もどそうとする力の強さ}\div\text{重さ}}$」の形です。弦は $\sqrt{F_{\mathrm T}/\mu}$、棒は $\sqrt{E/\rho}$、流体は $\sqrt{K/\rho}$。かたい（もどす力が強い）ほど速く、重いほどおそい。波-1 のばねとおもりの $\omega=\sqrt{k/m}$（この章の書き方では $\sqrt{k_{\mathrm s}/m}$）と同じ形です。
\end{phyomi}

\phboxhead{浅い水の波：津波}
水面の波では、もとにもどそうとする力は重力です。ここでは、波長が水の深さ $h$ よりずっと長い波（\textbf{浅い水の波}）を考えます。津波は、波長が数十〜数百 km もあるので、深さ\qty{4}{km}の海でも「浅い水の波」です。

\begin{bkfigure}
\begin{tikzpicture}[x=1mm, y=1mm]
  \fill[bkBlue!14] plot[smooth, domain=0:116, samples=60] (\x, {20+5/(1+((\x-58)/14)^2)}) -- (116,0) -- (0,0) -- cycle;
  \draw[bkBlue!70!black, line width=1pt] plot[smooth, domain=0:116, samples=60] (\x, {20+5/(1+((\x-58)/14)^2)});
  \draw[bkGray, line width=0.4pt, dashed] (0,20) -- (116,20);
  \fill[bkGray!50] (0,-2.4) rectangle (116,0);
  \fill[bkOrange!25] (36,0) rectangle (48,22.3);
  \draw[bkInk, line width=0.5pt] (36,0) -- (36,24) (48,0) -- (48,24);
  \foreach \yy in {4,10,16} {\draw[bkRed!80!black, line width=0.9pt, -{Stealth[length=1.8mm]}] (30,\yy) -- (35,\yy);}
  \foreach \yy in {4,10,16} {\draw[bkRed!80!black, line width=0.9pt, -{Stealth[length=1.8mm]}] (49,\yy) -- (54,\yy);}
  \node[font=\sffamily\footnotesize, anchor=east] at (29,10) {$u$};
  \draw[bkInk, line width=0.5pt, {Stealth[length=1.6mm]}-{Stealth[length=1.6mm]}] (8,0) -- (8,20);
  \node[font=\sffamily\footnotesize, anchor=west] at (9,10) {深さ $h$};
  \draw[bkInk, line width=0.5pt, {Stealth[length=1.4mm]}-{Stealth[length=1.4mm]}] (58,20) -- (58,25);
  \node[font=\sffamily\footnotesize, anchor=south west] at (59,25.3) {高さ $\eta$};
  \node[font=\sffamily\footnotesize, anchor=north] at (42,-3) {$x$ から $x+\Delta x$};
  \node[anchor=west, font=\small, text width=36mm, align=left] at (82,9) {水平の速さ $u$ は、底から水面まで同じ（波長≫深さ）};
\end{tikzpicture}
\bkfigcaption{浅い水の波。水面が $\eta$ だけ上がった所では、水の柱が少し高い。色の所（$x$ から $x+\Delta x$）に出入りする水と、両側の水圧の差を考える。}\label{ph:波-5:fig-asai}
\end{bkfigure}

\noindent 水面の高さ（静かなときからの上がり）を $\eta$（イータ）、水が横へ動く速さを $u$ とし、底から水面まで $u$ は同じとします。奥行き1 m あたりで考えます。
\begin{enumerate}[label={\bkcircled{\arabic*}}, leftmargin=2.4\zw, itemsep=2pt, topsep=2pt]
\item \textbf{水の量}：色の所に、左から1秒に $hu(x)$、右へ $hu(x+\Delta x)$ の水が流れます。その差だけ水面が上がるので、$\Delta x\,\partial\eta/\partial t=-h\,(\partial u/\partial x)\,\Delta x$、つまり
\[\frac{\partial\eta}{\partial t}=-h\,\frac{\partial u}{\partial x}\]
\item \textbf{運動方程式}：水圧は、水面からの深さに比例します（力-12）。左の水面が高いと、どの深さでも左の水圧が $\rho g\{\eta(x)-\eta(x+\Delta x)\}$ だけ大きいので、色の所は右へ押されます。高さ $h$ の柱全体では $-\rho gh\,(\partial\eta/\partial x)\,\Delta x$ の力です。質量は $\rho h\,\Delta x$ なので、
\[\frac{\partial u}{\partial t}=-g\,\frac{\partial\eta}{\partial x}\]
\end{enumerate}
①を $t$ で微分し、②を入れると、
\[\frac{\partial^2\eta}{\partial t^2}=-h\,\frac{\partial}{\partial x}\Bigl(\frac{\partial u}{\partial t}\Bigr)=gh\,\frac{\partial^2\eta}{\partial x^2}\]
で、また波動方程式です。だから浅い水の波の速さは
\pheq{c=\sqrt{gh}}{m/s ＝ $\sqrt{\text{m/s}^2\times\text{m}}$}
です。\bkidx{あさいみずのなみ}{浅い水の波}深さ\qty{4000}{m}の海では約\qty{200}{m/s}（時速約700 km、ジェット旅客機ほど）、深さ\qty{10}{m}では約\qty{10}{m/s}です（\bkreidairef{10}）。

\begin{phbaisuru}[4倍]{深さ $h$}
$c=\sqrt{gh}$ なので、速さは2倍です。反対に、津波が深い海から浅い海岸へ近づくと、どんどんおそくなります。
\par\smallskip\noindent\begin{tikzpicture}
\begin{axis}[bkplot, width=104mm, height=44mm, xmin=0, xmax=4400, ymin=0, ymax=230, xtick={1000,2000,3000,4000}, ytick={50,100,150,200}, xlabel={$h$（m）}, ylabel={$c$（m/s）}, scaled x ticks=false]
  \addplot[bkcurve, domain=0:4300, samples=200] {sqrt(9.8*x)};
  \fill[bkRed] (axis cs:1000,99) circle (1.8pt) (axis cs:4000,198) circle (1.8pt);
  \draw[bkGray, line width=0.4pt, dashed] (axis cs:1000,0) -- (axis cs:1000,99) (axis cs:4000,0) -- (axis cs:4000,198);
  \node[font=\sffamily\footnotesize, anchor=north west, fill=white, inner sep=1pt] at (axis cs:2000,120) {深さ4倍で速さ2倍};
  \node[font=\sffamily\footnotesize, anchor=south east, fill=white, inner sep=1pt] at (axis cs:960,103) {\qty{99}{m/s}};
  \node[font=\sffamily\footnotesize, anchor=south east, fill=white, inner sep=1pt] at (axis cs:3960,202) {\qty{198}{m/s}};
\end{axis}
\end{tikzpicture}
\end{phbaisuru}

\begin{phmijika}[津波は、なぜ海岸で高くなる？]
津波の先がおそくなると、後ろから来る部分が追いついてつまってきます。波長が短くなり、運んできたエネルギーが短い所に集まるので、波が高くなります。浅い水の波の理論では、高さは深さの4乗根に反比例し（グリーンの法則）、深さ\qty{4000}{m}から\qty{10}{m}まで来ると、約4.5倍の高さになります。沖では数十 cm でほとんど気づかない津波が、海岸では何 m にもなるのはこのためです。また、太平洋の向こうで起きた地震の津波も、ジェット機ほどの速さでやってきます。1960年に南米のチリで起きた大地震の津波は、約1日かけて太平洋をわたり、日本の海岸に大きな被害を出しました。
\end{phmijika}

\begin{phflow}[媒質の運動方程式から波の速さへ]
\phflowrow{$\rho\dfrac{\partial^2u}{\partial t^2}=E\dfrac{\partial^2u}{\partial x^2}$}{細い棒\\ひずみが小さい}{$c=\sqrt{\dfrac{E}{\rho}}$}
\phflowrow{$\rho\dfrac{\partial^2\xi}{\partial t^2}=K\dfrac{\partial^2\xi}{\partial x^2}$}{気体を、すばやく\\押す（$K_{\mathrm s}=\gamma p$）}{$c=\sqrt{\dfrac{\gamma p}{\rho}}$}
\phflowrow{$\dfrac{\partial\eta}{\partial t}=-h\dfrac{\partial u}{\partial x}$\\$\dfrac{\partial u}{\partial t}=-g\dfrac{\partial\eta}{\partial x}$}{波長≫深さ\\高さ≪深さ}{$c=\sqrt{gh}$}
\end{phflow}

\begin{phmodel}[棒・流体・浅い水の波]
\phignore{棒が細くなったり太くなったりすること（ポアソン比）。流体のねばり（粘性）。水底のでこぼこ。}
\phassume{ひずみ・圧力の変化・水面の高さは小さい（式が線形）。棒は細い（太さは波長よりずっと小さい）。水の波では、波長は深さよりずっと長い。}
\phlimit{太い棒や広がった固体（下のせいかくに言うと）。波長が深さと同じくらいか短い水の波（速さが波長で変わる。⑦）。海岸で波がくだけるほど高いとき（⑥）。}
\end{phmodel}

\begin{seikaku}
棒の縦波の $\sqrt{E/\rho}$ は、細い棒だけの値です。広がった固体（地面や厚い鉄の板）の中の縦波は、押された部分が横へふくらめないので、鋼では約\qty{5900}{m/s}と、棒より速くなります。固体の中の横波は、ずれの弾性（横弾性係数 $G$、力-8 ④）で伝わり、速さは $\sqrt{G/\rho}$ で、鋼では約\qty{3200}{m/s}です。地震の P 波が S 波より速いのも、同じ理由です。また、鋼の管に入った水の中の縦波（力-14 の水撃）は、水が押されると管がふくらむので、水だけのとき（約\qty{1500}{m/s}）よりおそく、\qtyrange{1000}{1300}{m/s}くらいになります。
\end{seikaku}

\begin{bkcenter}
\begin{bktable}{colspec={Q[l,m,wd=44mm] Q[c,m,wd=30mm] Q[c,m,wd=26mm] Q[c,m,wd=12mm]}, rowsep=1.2pt, row{1}={bg=bkTheme!10, font=\sffamily\small}}
媒質・波 & 式 & 速さ〔m/s〕 & 節 \\
ギターの A弦（110 Hz） & $\sqrt{F_{\mathrm T}/\mu}$ & 143 & ④ \\
空気（\qty{20}{\celsius}） & $\sqrt{\gamma p/\rho}$ & 343 & ⑤ \\
水 & $\sqrt{K/\rho}$ & 約1500 & ⑤ \\
鋼の棒（縦波） & $\sqrt{E/\rho}$ & 約5100 & ⑤ \\
広がった鋼（縦波・横波） & ― & 約5900・約3200 & ⑤ \\
津波（深さ\qty{4000}{m}） & $\sqrt{gh}$ & 約200 & ⑤ \\
\end{bktable}
\bktabcaption{いろいろな波の速さ。もどそうとする力が強く、軽い媒質ほど速い。}\label{ph:波-5:tab-hayasa}
\end{bkcenter}

\bksubsection{6}{重ね合わせ}
\label{ph:波-5:sec6}
\label{ph:波-5:kasane}向かい合った2つの波（パルス）がぶつかると、どうなるでしょう。

$y_1$ と $y_2$ が、どちらも波動方程式をみたすとします。2つをたした $y=y_1+y_2$ を入れると、
\[\frac{\partial^2(y_1+y_2)}{\partial t^2}=\frac{\partial^2y_1}{\partial t^2}+\frac{\partial^2y_2}{\partial t^2}=c^2\frac{\partial^2y_1}{\partial x^2}+c^2\frac{\partial^2y_2}{\partial x^2}=c^2\frac{\partial^2(y_1+y_2)}{\partial x^2}\]
で、和もまた波動方程式をみたします。微分は、和を項ごとに分けて計算できるからです（このような式を\textbf{線形}の式といいます。6-1）。だから、2つの波が同じ所に来たときの媒質のずれは、それぞれの波のずれの\textbf{和}になります。これを\textbf{重ね合わせの原理}といいます。\bkidx{かさねあわせのげんり}{重ね合わせの原理}

\begin{bkfigure}
\begin{tikzpicture}[x=1mm, y=1mm, declare function={rxft(\x,\c)=max(0,7-0.7*abs(\x-\c));}]
  \foreach \row/\L/\R/\lab in {0/20/80/{はじめ}, 1/43/57/{重なり始め}, 2/50/50/{ぴったり重なる}, 3/80/20/{通りぬけた後}} {%
    \begin{scope}[shift={(0,-\row*17)}]
    \draw[bkGray, line width=0.4pt] (0,0) -- (58,0);
    \draw[bkGray, line width=0.4pt] (64,0) -- (122,0);
    \draw[bkBlue!50, line width=0.6pt, dash pattern=on 2pt off 1pt, samples at={0,0.5,...,100}] plot ({0.58*\x}, {rxft(\x,\L)});
    \draw[bkOrange!70, line width=0.6pt, dash pattern=on 2pt off 1pt, samples at={0,0.5,...,100}] plot ({0.58*\x}, {rxft(\x,\R)});
    \draw[bkInk, line width=1.2pt, samples at={0,0.5,...,100}] plot ({0.58*\x}, {rxft(\x,\L)+rxft(\x,\R)});
    \draw[bkBlue!50, line width=0.6pt, dash pattern=on 2pt off 1pt, samples at={0,0.5,...,100}] plot ({64+0.58*\x}, {rxft(\x,\L)});
    \draw[bkOrange!70, line width=0.6pt, dash pattern=on 2pt off 1pt, samples at={0,0.5,...,100}] plot ({64+0.58*\x}, {-rxft(\x,\R)});
    \draw[bkInk, line width=1.2pt, samples at={0,0.5,...,100}] plot ({64+0.58*\x}, {rxft(\x,\L)-rxft(\x,\R)});
    \node[font=\sffamily\footnotesize, anchor=south west, fill=white, inner sep=0.5pt] at (0,8.5) {\lab};
    \end{scope}}
  \node[font=\sffamily\footnotesize, anchor=south] at (29,15) {（あ）山と山};
  \node[font=\sffamily\footnotesize, anchor=south] at (93,15) {（い）山と谷};
\end{tikzpicture}
\bkfigcaption{2つのパルスの重ね合わせ。破線はそれぞれの波、太い実線はその和（媒質の本当の形）。（あ）山と山がぴったり重なると高さは2倍、（い）山と谷では一瞬まっすぐになる。どちらも、通りぬけた後はもとの形にもどる。}\label{ph:波-5:fig-kasane}
\end{bkfigure}

\noindent（い）で、山と谷がぴったり重なった瞬間、弦はまっすぐになります。でも、波が消えたのではありません。弦の各部分は、そのとき\textbf{動いて}いて（速さは0ではない）、次の瞬間には、また山と谷に分かれて、それぞれ進んでいきます。2つの波は、ぶつかってもこわれず、通りぬけた後は、何事もなかったようにもとの形で進みます。これを\textbf{波の独立性}といいます。\bkidx{なみのどくりつせい}{波の独立性}\bkreidairef{5}で、重なっている途中の形を作図します。

\begin{phmijika}[いろいろな音が同時に聞こえる]
教室でたくさんの人が同時に話していても、1人1人の声を聞き分けられます。空気の中で、たくさんの音の波が重なり合っても、それぞれの波はこわれずに伝わってくるからです。ノイズキャンセリングのヘッドホンは、マイクで拾った周りの騒音と、ちょうど山と谷が逆になる音（逆位相の音）をスピーカーから出して重ね、耳の所で打ち消しています（くわしくは波-8）。
\end{phmijika}

\begin{seikaku}
重ね合わせの原理は、波動方程式が線形である（$y$ の2乗などがない）ことから出てきます。④⑤で「かたむきが小さい」「圧力の変化が小さい」と仮定したからこそ、線形の式になったのでした。ゆれが大きくなると、この仮定がくずれます。たとえば、海岸に近づいて高くなった波は、山の所の水が深い（$h+\eta$）ので、谷より速く進み（$\sqrt{gh}$）、前のめりになってくだけます。とても大きな音や爆発では、同じ理由で、圧力の高い所が追いついて急な段（衝撃波、熱-4 ⑨）ができます。このような波は、重ね合わせでは求められません（この先の学び）。
\end{seikaku}

\bksubsection{7}{ばねとおもりの列から波へ：分散関係と群速度}
\label{ph:波-5:sec7}
波-3 では、ばねでつながった2つ・3つのおもりのゆれを調べました。おもりの数を限りなく増やしたら、どうなるでしょう。

\begin{bkfigure}
\begin{tikzpicture}[x=1mm, y=1mm]
  \foreach \n/\lab in {0/{},1/{$n-1$},2/{$n$},3/{$n+1$},4/{}} {%
    \rxfmass{(10+\n*25,0)}
    \node[font=\sffamily\footnotesize, anchor=north] at (10+\n*25,-3.6) {\lab};}
  \foreach \n in {0,...,3} {\rxfspring{(12.4+\n*25,0)}{(32.6+\n*25,0)}}
  \draw[bkInk!60, line width=0.6pt, dotted] (-6,0) -- (7.6,0) (112.4,0) -- (124,0);
  \node[font=\sffamily\footnotesize, anchor=south] at (47.5,2.6) {$k_{\mathrm s}$};
  \node[font=\sffamily\footnotesize, anchor=south] at (72.5,2.6) {$k_{\mathrm s}$};
  \node[font=\sffamily\footnotesize, anchor=south] at (60,2.8) {$m$};
  \draw[bkInk, line width=0.5pt, {Stealth[length=1.6mm]}-{Stealth[length=1.6mm]}] (35,-9) -- (60,-9);
  \node[font=\sffamily\footnotesize, anchor=north] at (47.5,-9.4) {$a$};
  \draw[bkRed!80!black, line width=1pt, -{Stealth[length=1.8mm]}] (60,-14) -- (66,-14);
  \node[font=\sffamily\footnotesize, anchor=west, text=bkRed!80!black] at (66.5,-14) {$u_n$（右へのずれ）};
\end{tikzpicture}
\bkfigcaption{ばね（ばね定数 $k_{\mathrm s}$）とおもり（質量 $m$）が、間かく $a$ で限りなく並んだ列。$n$ 番めのおもりが、つり合いの位置から右へずれた量を $u_n$ とする。}\label{ph:波-5:fig-retsu}
\end{bkfigure}

\noindent この章では、波数を $k$ と書くので、ばね定数は $k_{\mathrm s}$ と書きます。$n$ 番めのおもりは、右のばねから $k_{\mathrm s}(u_{n+1}-u_n)$、左のばねから $-k_{\mathrm s}(u_n-u_{n-1})$ の力を受けます（右のおもりが自分より右へずれると、右のばねが伸びて右へ引く）。運動方程式は
\pheq{m\,\frac{\dd^2u_n}{\dd t^2}=k_{\mathrm s}\bigl(u_{n+1}-2u_n+u_{n-1}\bigr)}{N}
です。8-13 の「となりとの差の差」と同じ形です。

\phboxhead{長い波なら、波動方程式}
波長がおもりの間かく $a$ よりずっと長いと、となりどうしのずれは少ししか違いません。おもりの位置を $x=na$ とし、$u_n$ を、なめらかな関数 $u(x,\ t)$ の $x=na$ での値と考えます。5-17 の展開で、
\[u_{n\pm1}=u(x\pm a)\approx u\pm a\,\frac{\partial u}{\partial x}+\frac{a^2}{2}\,\frac{\partial^2u}{\partial x^2}\]
なので、$u_{n+1}-2u_n+u_{n-1}\approx a^2\,\partial^2u/\partial x^2$ です（8-12④の差の差）。運動方程式は
\[\frac{\partial^2u}{\partial t^2}=\frac{k_{\mathrm s}a^2}{m}\,\frac{\partial^2u}{\partial x^2}\qquad\text{だから}\qquad c_0=a\sqrt{\frac{k_{\mathrm s}}{m}}\]
の波動方程式になります。棒を、長さ $a$ の小さな部分に分けて、1つ1つをおもりとばねと考えると、$m=\rho Sa$、$k_{\mathrm s}=ES/a$（力-8 ③の棒のばね定数）なので、
\[c_0=a\sqrt{\frac{ES/a}{\rho Sa}}=\sqrt{\frac{E}{\rho}}\]
で、⑤の棒の波の速さにぴったりもどります。\textbf{連続した棒は、限りなく細かく分けたばねとおもりの列}なのです。力-16 で、車の部品を小さな要素（ばねとおもり）に分けて計算したのは、この考えです。

\begin{phflow}[ばねとおもりの列]
\phflowrow{$m\dfrac{\dd^2u_n}{\dd t^2}$\\$=k_{\mathrm s}(u_{n+1}-2u_n+u_{n-1})$}{波長≫間かく $a$\\（長い波）}{$c_0=a\sqrt{\dfrac{k_{\mathrm s}}{m}}$}
\phflowrow{$c_0=a\sqrt{\dfrac{k_{\mathrm s}}{m}}$}{棒を分けた\\$m=\rho Sa$、$k_{\mathrm s}=ES/a$}{$c=\sqrt{\dfrac{E}{\rho}}$}
\end{phflow}

\phboxhead{短い波まで：分散関係}
では、波長が短くて、$a$ と同じくらいのときはどうでしょう。正弦波 $u_n=A\cos(\omega t-kna)$ を、そのまま運動方程式に入れてみます。$\theta=\omega t-kna$ と書くと、$u_{n\pm1}=A\cos(\theta\mp ka)$ です。
\begin{bkhenkei}[0.56]
u_{n+1}+u_{n-1}=A\{\cos(\theta-ka)+\cos(\theta+ka)\} & となり2つ \\
=2A\cos\theta\cos ka=2u_n\cos ka & 加法定理（4-9②）で $\sin$ の項が消える \\
-m\omega^2u_n=k_{\mathrm s}(2\cos ka-2)\,u_n & 運動方程式に入れる（$\ddot u_n=-\omega^2u_n$） \\
\omega^2=\frac{2k_{\mathrm s}}{m}(1-\cos ka)=\frac{4k_{\mathrm s}}{m}\sin^2\frac{ka}{2} & 2倍角（4-9④）$1-\cos2\alpha=2\sin^2\alpha$ \\
\end{bkhenkei}
だから
\pheq{\omega=2\sqrt{\frac{k_{\mathrm s}}{m}}\,\Bigl|\sin\frac{ka}{2}\Bigr|}{rad/s ＝ $\sqrt{\text{N/m}\div\text{kg}}$}
\label{ph:波-5:bunsan}です。どの $n$ でも同じ式になったので、この $\omega$ と $k$ の組なら、正弦波は列をそのまま進めます。$\omega$ と $k$ の関係を表すこの式を、\textbf{分散関係}（ぶんさんかんけい）といいます。\bkidx{ぶんさんかんけい}{分散関係}波動方程式の分散関係は $\omega=ck$（④）でした。

\begin{bkfigure}
\begin{tikzpicture}
\begin{axis}[bkplot, width=112mm, height=56mm, xmin=0, xmax=3.6, ymin=0, ymax=2.5, xtick={0.7854,1.5708,3.1416}, xticklabels={$\pi/4$,$\pi/2$,$\pi$}, ytick={1,1.414,2}, yticklabels={1,$\sqrt2$,2}, xlabel={$ka$}, ylabel={$\omega\div\sqrt{k_{\mathrm s}/m}$}]
  \addplot[bkGray, line width=0.8pt, dashed, domain=0:2.45] {x};
  \addplot[bkcurve, domain=0:3.1416, samples=100] {2*sin(deg(x/2))};
  \addplot[bkOrange!85!black, line width=0.9pt, dash dot, domain=0.6:2.6] {1.4142+0.7071*(x-1.5708)};
  \addplot[bkRed!70!black, line width=0.9pt, dotted, domain=0:1.5708] {0.9003*x};
  \fill[bkRed] (axis cs:1.5708,1.4142) circle (2pt);
  \fill[bkInk] (axis cs:3.1416,2) circle (2pt);
  \node[font=\sffamily\footnotesize, anchor=south east, fill=white, inner sep=1pt] at (axis cs:2.3,2.35) {$\omega=c_0k$（波動方程式）};
  \node[font=\sffamily\footnotesize, anchor=north west, fill=white, inner sep=1pt] at (axis cs:1.75,1.2) {傾き＝群速度（接線）};
  \node[font=\sffamily\footnotesize, anchor=north west, fill=white, inner sep=1pt] at (axis cs:0.45,0.3) {傾き＝位相速度（原点から）};
  \node[font=\sffamily\footnotesize, anchor=north, fill=white, inner sep=1pt] at (axis cs:3.1416,1.9) {$\lambda=2a$};
\end{axis}
\end{tikzpicture}
\bkfigcaption{ばねとおもりの列の分散関係（太い実線）。破線は波動方程式（$\omega=c_0k$）。$ka=\pi/2$（$\lambda=4a$）の点で、原点からの点線の傾きが位相速度、曲線の接線（一点鎖線）の傾きが群速度。}\label{ph:波-5:fig-bunsan}
\end{bkfigure}

\begin{phyomi}[$\omega=2\sqrt{k_{\mathrm s}/m}\,|\sin(ka/2)|$]
波長が長い（$k$ が小さい）うちは、$\sin(ka/2)\approx ka/2$ で $\omega\approx c_0k$、つまり波動方程式と同じです。波長が短くなるほど、曲線は直線より下へ曲がり、$\lambda=2a$（$ka=\pi$）で $\omega$ がいちばん大きくなります。このとき、となりどうしのおもりは、いつも\textbf{反対向き}に動いています。
\end{phyomi}

\phboxhead{位相速度：短い波ほどおそい}
正弦波の山が進む速さは、$\omega t-kx$ が一定の点の速さ $\omega/k$ です。これを\textbf{位相速度}（いそうそくど）$c_{\mathrm p}$ といいます。\bkidx{いそうそくど}{位相速度}波動方程式なら $c_{\mathrm p}=c$ で、どの波長でも同じでした。ばねとおもりの列では、
\[c_{\mathrm p}=\frac{\omega}{k}=c_0\times\frac{\sin(ka/2)}{ka/2}\]
で、短い波ほどおそくなります。$a=\qty{0.10}{m}$、$\sqrt{k_{\mathrm s}/m}=\qty{20}{rad/s}$（$c_0=\qty{2.0}{m/s}$）の列で計算すると、表\ref{ph:波-5:tab-bunsan}のようになります。
\begin{bkcenter}
\begin{bktable}{colspec={Q[l,m,wd=34mm] *{5}{Q[c,m,wd=14mm]}}, rowsep=1.2pt, row{1}={bg=bkTheme!10, font=\sffamily\small}}
波長 $\lambda$ & $100a$ & $10a$ & $6a$ & $4a$ & $2a$ \\
位相速度 $c_{\mathrm p}$〔m/s〕 & 2.00 & 1.97 & 1.91 & 1.80 & 1.27 \\
群速度 $c_{\mathrm g}$〔m/s〕 & 2.00 & 1.90 & 1.73 & 1.41 & 0 \\
\end{bktable}
\bktabcaption{ばねとおもりの列の位相速度と群速度（$c_0=\qty{2.0}{m/s}$）。波長がおもりの間かくの10倍より長ければ、どちらもほぼ $c_0$。}\label{ph:波-5:tab-bunsan}
\end{bkcenter}
波長で速さが変わることを、波の\textbf{分散}（ぶんさん）といいます。\bkidx{ぶんさん}{分散（波の）}いろいろな波長をふくむパルスは、波長ごとに違う速さで進むので、進むうちに形がくずれて広がります。分散のない波動方程式の波（③）だけが、形を保って進むのです。

\phboxhead{うなりのつつみの速さ：群速度}
\label{ph:波-5:gun}波長の少し違う2つの正弦波を重ねてみます。4-9⑦の積と和の式
\[\cos\alpha+\cos\beta=2\cos\frac{\alpha-\beta}{2}\cos\frac{\alpha+\beta}{2}\]
を使うと、
\[\cos(\omega_1t-k_1x)+\cos(\omega_2t-k_2x)\]
\[=2\cos\Bigl(\frac{\Delta\omega}{2}t-\frac{\Delta k}{2}x\Bigr)\cos(\bar\omega t-\bar kx)\]
です（$\Delta\omega=\omega_1-\omega_2$、$\Delta k=k_1-k_2$、$\bar\omega$・$\bar k$ は2つの平均）。右の $\cos(\bar\omega t-\bar kx)$ は細かい波で、速さ $\bar\omega/\bar k$（位相速度）で進みます。左の $\cos(\cdots)$ は、細かい波の大きさをゆっくり変える「つつみ」（うなりのかたまり。4-9⑦）で、$\Delta\omega\,t/2-\Delta k\,x/2$ が一定の所、つまり速さ $\Delta\omega/\Delta k$ で進みます。2つの波長を近づけた先の値
\pheq{c_{\mathrm g}=\frac{\dd\omega}{\dd k}}{m/s ＝ $\dfrac{\text{rad/s}}{\text{rad/m}}$}
を\textbf{群速度}（ぐんそくど）といいます。\bkidx{ぐんそくど}{群速度}分散関係のグラフの\textbf{接線の傾き}です（図\ref{ph:波-5:fig-bunsan}）。

\begin{bkfigure}
\begin{tikzpicture}
\begin{axis}[bkplot, width=120mm, height=32mm, xmin=-1.0, xmax=2.2, ymin=-2.4, ymax=2.9, xtick={-0.5,0,0.5,1,1.5,2}, ytick=\empty, axis x line*=bottom, axis y line=none, xlabel={$x$（m）}]
  \addplot[bkcurve, line width=1pt, domain=-1:2.2, samples=500] {cos(deg(-14.28*x))+cos(deg(-17.4533*x))};
  \addplot[bkOrange!85!black, line width=0.8pt, dash pattern=on 3pt off 1.5pt, domain=-1:2.2, samples=200] {2*abs(cos(deg(1.5867*x)))};
  \fill[bkRed] (axis cs:0,2) circle (2pt);
  \node[font=\footnotesize, anchor=south, text=bkOrange!85!black] at (axis cs:0,2.05) {▼};
  \node[font=\sffamily\footnotesize, anchor=west, fill=white, inner sep=1pt] at (axis cs:-1,2.6) {$t=0$};
\end{axis}
\end{tikzpicture}\par\vspace{2mm}
\begin{tikzpicture}
\begin{axis}[bkplot, width=120mm, height=32mm, xmin=-1.0, xmax=2.2, ymin=-2.4, ymax=2.9, xtick={-0.5,0,0.5,1,1.5,2}, ytick=\empty, axis x line*=bottom, axis y line=none, xlabel={$x$（m）}]
  \addplot[bkcurve, line width=1pt, domain=-1:2.2, samples=500] {cos(deg(26.1944*0.5-14.28*x))+cos(deg(30.6418*0.5-17.4533*x))};
  \addplot[bkOrange!85!black, line width=0.8pt, dash pattern=on 3pt off 1.5pt, domain=-1:2.2, samples=200] {2*abs(cos(deg(-2.2237*0.5+1.5867*x)))};
  \fill[bkRed] (axis cs:0.8955,1.88) circle (2pt);
  \node[font=\footnotesize, anchor=south, text=bkOrange!85!black] at (axis cs:0.7007,2.05) {▼};
  \draw[bkGray, line width=0.4pt, dashed] (axis cs:0,-2.4) -- (axis cs:0,2.9);
  \node[font=\sffamily\footnotesize, anchor=west, fill=white, inner sep=1pt] at (axis cs:-1,2.6) {$t=\qty{0.5}{s}$};
\end{axis}
\end{tikzpicture}
\bkfigcaption{ばねとおもりの列（$c_0=\qty{2.0}{m/s}$）で、波長 $4.4a$ と $3.6a$ の波を重ねた形。破線がつつみ。▼はつつみの山、●は細かい波の1つの山。0.5秒で（下の図。破線は $x=0$）、●は約\qty{0.90}{m}（位相速度 約\qty{1.8}{m/s}）、▼は約\qty{0.70}{m}（群速度 約\qty{1.4}{m/s}）進む。}\label{ph:波-5:fig-gun}
\end{bkfigure}

\noindent ばねとおもりの列では、分散関係を $k$ で微分して（5-15）、
\[c_{\mathrm g}=\frac{\dd\omega}{\dd k}=a\sqrt{\frac{k_{\mathrm s}}{m}}\cos\frac{ka}{2}=c_0\cos\frac{ka}{2}\]
です。表\ref{ph:波-5:tab-bunsan}の下の行がこの値で、群速度は位相速度よりさらにおそく、$\lambda=2a$ では0になります。となりどうしが反対に動くゆれは、その場でふるえるだけで、先へ進まないのです。

波が運ぶエネルギーや、波にのせた信号（音の始まりや、言葉の区切り）は、細かい山ではなく、つつみといっしょに進みます。だから、\textbf{エネルギーや信号が進む速さは群速度}です（くわしくは波-7）。分散のない波（波動方程式）では、$\omega=ck$ の傾きはどこでも $c$ なので、$c_{\mathrm g}=c_{\mathrm p}=c$ です。

\begin{phmijika}[分散する波は、身の回りにたくさんある]
池に石を投げこむと、波の輪が広がっていくうちに、長い波が先に、短い波が後に並んでいきます。水の波は、波長が深さより短くなると、長い波ほど速い（分散する）からです。深い水の波では、群速度は位相速度の半分になります（波-7 ⑥）。ガラスの中の光も、色（波長）で速さが少し違い、プリズムで白い光が色に分かれます（波-12）。はりや板を伝わる曲げの波も、波長で速さが変わります（波-6 ⑨）。
\end{phmijika}

\phboxhead{いちばん速いゆれと、計算の刻み}
分散関係のグラフで、$\omega$ がいちばん大きいのは $\lambda=2a$ のときで、
\pheq{\omega_{\max}=2\sqrt{\frac{k_{\mathrm s}}{m}}=\frac{2c_0}{a}}{rad/s}
です。力-16 ⑦で「棒を小さなばねとおもりの列とみると、いちばん速いゆれは、$2c$ を要素の大きさでわった角振動数」と書いたのは、この式です。速度を先に進めるオイラー法（中心差分）でこわれない条件は $\omega\,\Delta t<2$ でした（力-16 ⑦、7-10⑨）。いちばん速いゆれでもこわれないようにすると、
\[\frac{2c_0}{a}\,\Delta t<2\qquad\text{つまり}\qquad\Delta t<\frac{a}{c_0}\]
で、力-16 ⑦の CFL 条件「刻みは、波が1つの要素をわたる時間より短く」が出てきます（\bkreidairef{11}）。

\begin{phbaisuru}[半分]{おもりの間かく $a$（同じ棒を細かく分ける）}
同じ棒なら $c_0=\sqrt{E/\rho}$ は変わらないので、$\omega_{\max}=2c_0/a$ は2倍になり、こわれない刻みの上限 $a/c_0$ は半分になります。細かく分けるほど、細かい波まで表せますが、そのぶん時間の刻みも小さくしなければなりません（力-16 ⑦）。
\end{phbaisuru}

\begin{seikaku}
分散関係の $k$ を、$\pi/a$ より大きくしても、新しいゆれは出てきません。たとえば $k$ と $2\pi/a-k$ は、おもりの位置（$x=na$）だけで見ると、まったく同じずれ方になるからです（間の点はないので、区別できない）。おもりの列が表せる波は、$\lambda=2a$ より長い波だけです。動画の車輪が逆に回って見えるのと同じしくみです（波-11）。
\end{seikaku}

\bksubsection{8}{表計算で波を動かす}
\label{ph:波-5:sec8}
8-13⑧で作った波の表計算を、⑦の目で見直しましょう。弦を間かく $\Delta x$ の点に分け、時刻を $\Delta t$ ごとに区切ります。$n$ 歩めの点 $i$ のずれを $y_i^{\,n}$ と書くと、1歩の式は
\[y_i^{\,n+1}=2y_i^{\,n}-y_i^{\,n-1}+s^2\bigl(y_{i+1}^{\,n}-2y_i^{\,n}+y_{i-1}^{\,n}\bigr),\qquad s=\frac{c\,\Delta t}{\Delta x}\]
でした（8-13⑧）。右辺の $y_{i+1}-2y_i+y_{i-1}$ は、⑦のばねとおもりの列の力そのものです。時間の進め方は、力-16 ⑦の「速度を先に進めるオイラー法」と同じ計算になります（速度 $(y^{\,n+1}-y^{\,n})/\Delta t$ を先に新しくして、位置を進めている）。つまり、波の表計算は、\textbf{弦をおもりの列とみて、1つ1つのおもりの運動方程式を解いている}のです。

\phboxhead{表計算の作り方}
A列を歩の番号、B列からV列を点 $i=0$〜20 とします。両はし（B列・V列）は固定して、いつも0です。2行めに、はじめの形（点9・10・11 を0.5・1・0.5、ほかは0）を入れます。手をはなした瞬間は止まっているので、3行め（1歩め）は
\[y^{\,1}_i=y^{\,0}_i+\frac{s^2}{2}\bigl(y_{i+1}^{\,0}-2y_i^{\,0}+y_{i-1}^{\,0}\bigr)\]
とします（8-13⑧で $s=1$ のときに「となり2つの平均」とした式）。$s$ の値をX1のセルに入れて、4行めのC列に
\begin{bkcenter}
\texttt{C4 = 2*C3 - C2 + \$X\$1\^{}2*(D3 - 2*C3 + B3)}
\end{bkcenter}
と入れ、右と下へコピーします。$s=1$ と $s=1.1$ で計算した結果が、表\ref{ph:波-5:tab-sheet}と図\ref{ph:波-5:fig-sheet}です。

\begin{bkcenter}
\begin{bktable}{colspec={Q[c,m,wd=13mm] Q[c,m,wd=7mm] *{9}{Q[c,m,wd=10.6mm]}}, colsep=1.5pt, rowsep=0.8pt, row{1}={bg=bkTheme!10, font=\sffamily\small}, cell{2}{1}={r=5}{c}, cell{7}{1}={r=5}{c}, hline{7}={0.4pt, bkGray}}
 & 歩 & 点6 & 点7 & 点8 & 点9 & 点10 & 点11 & 点12 & 点13 & 点14 \\
$s=1$ & 0 & 0 & 0 & 0 & 0.5 & 1 & 0.5 & 0 & 0 & 0 \\
 & 1 & 0 & 0 & 0.25 & 0.5 & 0.5 & 0.5 & 0.25 & 0 & 0 \\
 & 2 & 0 & 0.25 & 0.5 & 0.25 & 0 & 0.25 & 0.5 & 0.25 & 0 \\
 & 3 & 0.25 & 0.5 & 0.25 & 0 & 0 & 0 & 0.25 & 0.5 & 0.25 \\
 & 4 & 0.5 & 0.25 & 0 & 0 & 0 & 0 & 0 & 0.25 & 0.5 \\
$s=1.1$ & 0 & 0 & 0 & 0 & 0.5 & 1 & 0.5 & 0 & 0 & 0 \\
 & 1 & 0 & 0 & 0.30 & 0.50 & 0.40 & 0.50 & 0.30 & 0 & 0 \\
 & 2 & 0 & 0.37 & 0.48 & 0.13 & 0.04 & 0.13 & 0.48 & 0.37 & 0 \\
 & 3 & 0.44 & 0.43 & 0.10 & 0.08 & $-0.09$ & 0.08 & 0.10 & 0.43 & 0.44 \\
 & 4 & 0.33 & 0.12 & 0.08 & $-0.15$ & 0.18 & $-0.15$ & 0.08 & 0.12 & 0.33 \\
\end{bktable}
\bktabcaption{波の表計算（点6〜14 の所だけ。$s=1.1$ は小数第2位まで）。$s=1$ では、高さ0.5の山が2つに分かれ、1歩に1点ずつ左右へ進む。$s=1.1$ では、まん中に $+$ と $-$ が交互に並ぶぎざぎざが出始める。}\label{ph:波-5:tab-sheet}
\end{bkcenter}

\begin{bkfigure}
\begin{tikzpicture}
\begin{axis}[bkplot, name=s1, width=62mm, height=46mm, xmin=0, xmax=20.5, ymin=-1.1, ymax=1.15, xtick={5,10,15,20}, ytick={-1,0.5,1}, ylabel={$y$}, title={$s=1$}, title style={font=\sffamily\footnotesize}, axis x line*=bottom]
  \draw[bkInk, line width=0.5pt] (axis cs:0,0) -- (axis cs:20.5,0);
  \addplot[bkGray, line width=0.7pt, mark=square*, mark size=1.2pt] coordinates {(0,0) (8,0) (9,0.5) (10,1) (11,0.5) (12,0) (20,0)};
  \addplot[bkOrange!85!black, line width=0.9pt, mark=triangle*, mark size=1.6pt] coordinates {(0,0) (5,0) (6,0.25) (7,0.5) (8,0.25) (9,0) (11,0) (12,0.25) (13,0.5) (14,0.25) (15,0) (20,0)};
  \addplot[bkNavy, line width=1pt, mark=*, mark size=1.2pt] coordinates {(0,0) (2,0) (3,0.25) (4,0.5) (5,0.25) (6,0) (14,0) (15,0.25) (16,0.5) (17,0.25) (18,0) (20,0)};
\end{axis}
\begin{axis}[bkplot, at={(s1.east)}, anchor=west, xshift=12mm, width=62mm, height=46mm, xmin=0, xmax=20.5, ymin=-3.3, ymax=3.4, xtick={5,10,15,20}, ytick={-3,-1,1,3}, title={$s=1.1$}, title style={font=\sffamily\footnotesize}, axis x line*=bottom]
  \draw[bkInk, line width=0.5pt] (axis cs:0,0) -- (axis cs:20.5,0);
  \addplot[bkGray, line width=0.7pt, mark=square*, mark size=1.2pt] coordinates {(0,0) (8,0) (9,0.5) (10,1) (11,0.5) (12,0) (20,0)};
  \addplot[bkOrange!85!black, line width=0.9pt, mark=triangle*, mark size=1.6pt] coordinates {(0,0) (4,0) (5,0.536) (6,0.328) (7,0.115) (8,0.084) (9,-0.151) (10,0.176) (11,-0.151) (12,0.084) (13,0.115) (14,0.328) (15,0.536) (16,0) (20,0)};
  \addplot[bkNavy, line width=1pt, mark=*, mark size=1.2pt] coordinates {(0,0) (1,1.149) (2,-0.892) (3,1.557) (4,-1.560) (5,1.191) (6,-0.385) (7,-0.714) (8,1.841) (9,-2.683) (10,2.995) (11,-2.683) (12,1.841) (13,-0.714) (14,-0.385) (15,1.191) (16,-1.560) (17,1.557) (18,-0.892) (19,1.149) (20,0)};
\end{axis}
\end{tikzpicture}
\bkfigcaption{波の表計算（横の軸は点の番号 $i$）。左は $s=1$（■0歩、▲3歩、●6歩）、右は $s=1.1$（■0歩、▲4歩、●8歩。たての目もりが左の3倍）。$s=1.1$ では、ぎざぎざが1歩に約2.4倍ずつ育ち、10歩で高さ約13、20歩で約4万になる。}\label{ph:波-5:fig-sheet}
\end{bkfigure}

\phboxhead{なぜ $s\leqq1$ なのか}
こわれ方を見ると、$+$ と $-$ が1点ごとに入れかわる、いちばん細かいぎざぎざ（$\lambda=2\Delta x$）が育っています。これは、⑦の列のいちばん速いゆれ $\omega_{\max}=2c/\Delta x$ です。このゆれがこわれない条件 $\omega_{\max}\Delta t\leqq2$ は、
\[\frac{2c}{\Delta x}\,\Delta t\leqq2\qquad\text{つまり}\qquad s=\frac{c\,\Delta t}{\Delta x}\leqq1\]
で、8-13⑧の「1歩のあいだに、波がマス1つより遠くへ行ってはいけない」と同じ条件になります。8-13 では「となりのマスしか見ない計算は、2マス先へ一気に進む波を表せない」と説明しました。ここでは、同じ条件が「いちばん速いゆれを、刻みが追いきれるか」としても読めることがわかりました。

\begin{seikaku}
$s=1$ のとき、山は形をまったく変えずに進みました。これは1次元の波の表計算だけの、特別にうまくいく場合です。$s<1$（たとえば0.5）にすると、こわれはしませんが、山は進むうちに少し低く広がり、後ろに小さなさざ波を引きずります。⑦で見たように、おもりの列では短い波ほどおそい（分散する）からです。これを\textbf{数値分散}といい、本物の弦にはない、計算が作ったにせものの分散です。マスを細かくするほど小さくなります。
\end{seikaku}

\bksubsection{9}{複素数で書く波}
\label{ph:波-5:sec9}
6-10⑨で、交流を回る矢印（フェーザ）で書くと、計算が楽になることを学びました。波でも同じことができます。
\pheq{y=\mathrm{Re}\bigl[X\,e^{j(\omega t-kx)}\bigr]\qquad X=Ae^{j\varphi}}{m}
と書きます（Re は実部。6-9）。$X$ を\textbf{複素振幅}といい、大きさ $|X|=A$ が振幅、角 $\varphi$ が初期位相です。$e^{j(\omega t-kx)}$ を $t$ で微分すると $j\omega$ が、$x$ で微分すると $-jk$ がかかります（6-10⑩）。だから、微分が「かけ算」になります。

\phboxhead{波動方程式に入れる}
\[(j\omega)^2X\,e^{j(\omega t-kx)}=c^2(-jk)^2X\,e^{j(\omega t-kx)}\qquad\text{だから}\qquad-\omega^2=-c^2k^2,\quad\omega=ck\]
1行で、④の分散関係 $\omega=ck$ が出ました。

\phboxhead{ばねとおもりの列に入れる}
$u_n=\mathrm{Re}[X\,e^{j(\omega t-kna)}]$ とすると、となりの $u_{n\pm1}$ には $e^{\mp jka}$ がかかるだけです。
\[-m\omega^2=k_{\mathrm s}\bigl(e^{-jka}-2+e^{jka}\bigr)=k_{\mathrm s}(2\cos ka-2)\]
（6-10⑧：$e^{j\theta}+e^{-j\theta}=2\cos\theta$）。⑦で加法定理を使った計算が、ずっと短くなりました。専門の本では、波をほとんどいつもこの形で書きます。

\begin{seikaku}
波が進むうちに弱まっていく（吸収される）ときは、$k$ を複素数 $k=k_{\mathrm r}-j\alpha$ にすると表せます。$e^{-jkx}=e^{-\alpha x}e^{-jk_{\mathrm r}x}$ となり、振幅が $e^{-\alpha x}$ で小さくなっていく波です（波-7 ⑦）。また、空間を進む波は、$x$ のかわりにベクトルを使って $e^{j(\omega t-\vec k\cdot\vec r)}$ と書きます。$\vec k$ は、波の進む向きを向いた、長さ $2\pi/\lambda$ のベクトル（波数ベクトル）です。
\end{seikaku}

\begin{phmada}[ウェーブは、なぜ時計回りが多い？]
スタジアムのウェーブを、たくさんの試合のビデオで調べた研究があります（2002年、ファルカシュら）。大きなスタジアムでは、25〜35人くらいが同時に立ち上がるとウェーブが始まり、1秒に約20席の速さで進むことがわかりました。そして、調べたウェーブの多くが、上から見て時計回りに進んでいた、と報告されています。研究者たちは、人を「まわりの何人かが立ち上がると、自分も立ち上がり、しばらくは休む」という、山火事の燃え広がりや心臓の筋肉の興奮と同じ形のモデルで表し、速さや幅をよく再現しました。でも、なぜ時計回りが多いのかは、このモデルからは出てきません。研究者たちは、右ききの人が多いことと関係があるかもしれないと書いていますが、確かめられてはおらず、はっきりしていません。人の群れの動きを物理の式で表す研究は、今も続いています。
\end{phmada}

\begin{phsaki}{非線形の波・電磁波・物質の波}
この章の波は、どれも「小さなゆれ」で、重ね合わせができる線形の波でした。1834年、イギリスのラッセルは、運河で、形をくずさずに進む1つの水の盛り上がり（高さ約\qty{40}{cm}）を、馬で2 km ほど追いかけました。分散で広がろうとする力と、高い所ほど速いという非線形の効果がつり合ってできるこの波は、今では\textbf{ソリトン}とよばれ、光ファイバーの通信にも使われています（専門の本で）。また、真空でも伝わる光（電磁波）は、物理シリーズの『電磁気学』（電-15）で、電気と磁気の式から同じ形の波動方程式として導きます。電子のような小さな粒も波としてふるまうことは、物理シリーズの『原子』（原-3）で学びます。平面や空間に広がる波は、$\partial^2/\partial x^2$ を $\nabla^2$（8-12）にかえた式で表します（波-7・波-8）。
\end{phsaki}

\begin{senmonsho}[振動・音響の専門書の書き方]
\begin{bkhonbun}
一次元の波動方程式
\[\frac{\partial^2u}{\partial t^2}=c^2\,\frac{\partial^2u}{\partial x^2}\]
の一般解は $u=f(x-ct)+g(x+ct)$ である．調和波 $u=\mathrm{Re}[A\,e^{i(kx-\omega t)}]$ を代入すると分散関係 $\omega=\omega(k)$ が得られ，位相速度は $v_{\mathrm p}=\omega/k$，群速度は $v_{\mathrm g}=\dd\omega/\dd k$ で与えられる．
\end{bkhonbun}
\bkyaku{物理の専門書は、虚数単位を $i$ と書き、$e^{i(kx-\omega t)}$ のように $x$ と $t$ の順と符号を入れかえて書くことが多い（実部をとれば、この章の $\cos(\omega t-kx)$ と同じ波）。電気の本は $j$ を使い、$e^{j(\omega t-kx)}$ と書く。位相速度・群速度を $v_{\mathrm p}$・$v_{\mathrm g}$ と書く本も多い（この本の $c_{\mathrm p}$・$c_{\mathrm g}$）。「調和波」は正弦波のこと。}
\end{senmonsho}

\begin{taisetsu}
\def\tsl#1{\makebox[7.6\zw][l]{\sffamily\bfseries\small\color{\bkTT}#1}}
\tsl{波の速さ}$c=f\lambda=\dfrac{\lambda}{T}$。$c$ は媒質で、$f$ は波源で決まる\\[3pt]
\tsl{正弦波}$y=A\sin(\omega t-kx)$、$k=\dfrac{2\pi}{\lambda}$、$\omega=ck$\\[3pt]
\tsl{波動方程式}$\dfrac{\partial^2y}{\partial t^2}=c^2\dfrac{\partial^2y}{\partial x^2}$。$y_0(x-ct)$ は形を保って進む\\[3pt]
\tsl{波の速さの式}弦 $\sqrt{\dfrac{F_{\mathrm T}}{\mu}}$、棒 $\sqrt{\dfrac{E}{\rho}}$、流体 $\sqrt{\dfrac{K}{\rho}}$、浅い水 $\sqrt{gh}$\\[3pt]
\tsl{重ね合わせ}ずれは、それぞれの波のずれの和。通りぬけた後はもとの形\\[3pt]
\tsl{分散と群速度}おもりの列 $\omega=2\sqrt{\dfrac{k_{\mathrm s}}{m}}\Bigl|\sin\dfrac{ka}{2}\Bigr|$。$c_{\mathrm p}=\dfrac{\omega}{k}$、$c_{\mathrm g}=\dfrac{\dd\omega}{\dd k}$\\[3pt]
\tsl{計算の刻み}$s=\dfrac{c\,\Delta t}{\Delta x}\leqq1$（いちばん速いゆれ $\omega_{\max}=\dfrac{2c}{\Delta x}$）
\end{taisetsu}

\bkkigocards{\bkkigocard{$c$}{シー}{波の速さ〔m/s〕。この巻の第2・3部の意味（第1部では減衰係数、第4部では光速）。高校の教科書の $v$・$V$}{wave speed}{波-5}}
{\bkkigocard{$k$}{ケー}{波数 $2\pi/\lambda$〔rad/m〕。波-5 では、ばね定数を $k_{\mathrm s}$ と書いて分ける}{wavenumber}{波-5}}
\bkkigocards{\bkkigocard{$\lambda$}{ラムダ}{波長〔m〕。固有値の $\lambda$（波-3）とは別の量}{wavelength}{波-5}}
{\bkkigocard{$\mu$}{ミュー}{線密度（弦の 1 m あたりの質量）〔kg/m〕。数学 8-13 と同じ。摩擦係数の $\mu$（力-2）とは別}{linear density}{波-5}}
\bkkigocards{\bkkigocard{$c_{\mathrm p}$・$c_{\mathrm g}$}{シー・ピー、シー・ジー}{位相速度 $\omega/k$ と群速度 $\dd\omega/\dd k$〔m/s〕。専門書は $v_{\mathrm p}$・$v_{\mathrm g}$}{phase / group velocity}{波-5}}
{\bkkigocard{$F_{\mathrm T}$}{エフ・ティー}{張力〔N〕。数学 8-13 は $F$ と書く。周期 $T$ とぶつけないための添字}{tension}{波-5}}

\begin{kotoba}
\small
\bkword{媒質}{波を伝える物（弦・水・空気・地面など）。媒質そのものは、波といっしょに遠くへは運ばれない。}
\bkword{横波・縦波}{媒質のゆれる向きが、波の進む向きに垂直な波が横波、平行な波が縦波（疎密波）。}
\bkword{波長・波数}{となりの山までの長さ $\lambda$。1 m あたりの角の変わり方 $k=2\pi/\lambda$。}
\bkword{波動方程式}{媒質の小さな部分の運動方程式から作った、波を表す式。}
\bkword{分散関係}{波の角振動数 $\omega$ と波数 $k$ の関係。$\omega=ck$ でないと、波長で速さが変わる（分散）。}
\bkword{位相速度・群速度}{山の進む速さ $\omega/k$ と、つつみ（エネルギー・信号）の進む速さ $\dd\omega/\dd k$。}
\end{kotoba}

%% ======================== 例題 ========================
\Needspace{10\baselineskip}\bksection{reidai}{例題でたしかめよう}
\label{ph:波-5:ex1}
\begin{phreidai}[スタジアムのウェーブ]
ウェーブは、1秒に約20席の速さで進み、その速さは約\qty{12}{m/s}です。（1）席1つの幅は何 m ですか。（2）立ち上がっている部分の幅が\qty{9}{m}のとき、1人が立っている時間は何秒ですか。
\phbamen
観客の「立ち上がる」という形が、客席にそって進む場面です（①）。
\phmushi
席の幅はどこも同じで、ウェーブは一定の速さで、形を変えずに進むとします。
\phzu
速さ $c=\qty{12}{m/s}$、1秒に $N=20$ 席、席の幅 $d$、立っている部分の幅 $w=\qty{9}{m}$、立っている時間 $\tau$。
\phriyu
1秒に進む距離（速さ）は、1秒に進む席の数×席の幅です。また、1人の所を、立っている部分（幅 $w$）が通りすぎるあいだ、その人は立っています。
\phtateru
\[c=Nd,\qquad\tau=\frac{w}{c}\]
\phtoku
（1）\[d=\frac{c}{N}=\frac{12}{20}=\qty{0.6}{m}\]
（2）\[\tau=\frac{9}{12}=\qty{0.75}{s}\]
\phtashikame
単位：（m/s）÷（席/s）＝ m/席、m÷（m/s）＝ s。\\
大きさ：席の幅\qty{0.6}{m}は、いすと少しのすき間の幅として、ありそうな大きさです。1秒足らず立ってすわるのも、実際のウェーブの動きに合います。
\kotae （1）約\qty{0.6}{m}、（2）約\qty{0.75}{s}
\end{phreidai}
\begin{jibun}
1周が\qty{540}{m}のスタジアムを、例題1のウェーブ（\qty{12}{m/s}）が1周するのに何秒かかりますか。その間に、ウェーブは何席分進みますか。
\bkjans{$540\div12=\qty{45}{s}$。$20\times45=900$ 席分}
\end{jibun}

\label{ph:波-5:ex2}
\begin{phreidai}[ラの音の波長]
時報や楽器の音合わせに使うラの音は、振動数\qty{440}{Hz}です。（1）\qty{20}{\celsius}の空気（音の速さ\qty{343}{m/s}）の中での波長を求めましょう。（2）同じ音が水（音の速さ\qty{1500}{m/s}）の中に伝わると、波長はいくらですか。
\phbamen
音さやスピーカーが空気をゆらし、その疎密波が伝わる場面です。
\phmushi
空気の温度はどこも同じで、音の速さは一定とします。
\phzu
$f=\qty{440}{Hz}$、空気中 $c_1=\qty{343}{m/s}$、水中 $c_2=\qty{1500}{m/s}$、波長 $\lambda_1$・$\lambda_2$。
\phriyu
1周期に1波長進むので $c=f\lambda$（②）。振動数は波源（音さ）が決めるので、水の中でも\qty{440}{Hz}のままです。
\phtateru
\[\lambda=\frac{c}{f}\]
\phtoku
（1）\[\lambda_1=\frac{343}{440}\approx\qty{0.78}{m}\]
（2）\[\lambda_2=\frac{1500}{440}\approx\qty{3.4}{m}\]
\phtashikame
単位：（m/s）÷ Hz ＝（m/s）×s ＝ m。\\
大きさ：空気の中のラの音の波長は、両手を広げたくらいです。水の中では、速さが約4.4倍なので、波長も約4.4倍になりました。
\kotae （1）約\qty{0.78}{m}、（2）約\qty{3.4}{m}
\end{phreidai}
\begin{jibun}
（1）人に聞こえるいちばん低い音（\qty{20}{Hz}）といちばん高い音（\qty{20}{kHz}）の、空気（\qty{343}{m/s}）の中での波長を求めましょう。（2）波長\qty{1.2}{m}、周期\qty{0.40}{s}の波の速さと振動数を求めましょう。
\bkjans{（1）$343\div20\approx\qty{17}{m}$、$343\div\num{20000}\approx\qty{0.017}{m}$（\qty{17}{mm}）。（2）$c=1.2\div0.40=\qty{3.0}{m/s}$、$f=1\div0.40=\qty{2.5}{Hz}$}
\end{jibun}

\label{ph:波-5:ex3}
\begin{phreidai}[正弦波の式を読む]
ぴんと張ったひもを伝わる波が、$y=0.02\sin(10\pi t-5\pi x)$（$x$・$y$ は m、$t$ は s）で表されます。（1）振幅・角振動数・波数・振動数・波長・波の速さを求めましょう。（2）$x=\qty{0.05}{m}$、$t=\qty{0.05}{s}$ のずれを求めましょう。（3）ひもの各点が上下に動く速さの最大値を求めましょう。
\phbamen
式で書かれた波から、波の要素と、ある場所・ある時刻のようすを読みとる場面です。
\phmushi
ひもは十分に長く、はね返ってくる波はないとします。
\phzu
$y=A\sin(\omega t-kx)$ とくらべる。$A$、$\omega$、$k$、$f$、$\lambda$、$c$。
\phriyu
③の正弦波の式です。$\omega=2\pi f$、$k=2\pi/\lambda$、$c=\omega/k$。媒質の速さは $\partial y/\partial t$ で、最大は $A\omega$（③）。
\phtateru
\[A=0.02,\quad\omega=10\pi,\quad k=5\pi,\qquad f=\frac{\omega}{2\pi},\quad\lambda=\frac{2\pi}{k},\quad c=\frac{\omega}{k}\]
\phtoku
（1）$A=\qty{0.02}{m}$、$\omega=10\pi\approx\qty{31}{rad/s}$、$k=5\pi\approx\qty{16}{rad/m}$。
\[f=\frac{10\pi}{2\pi}=\qty{5}{Hz},\qquad\lambda=\frac{2\pi}{5\pi}=\qty{0.4}{m},\qquad c=\frac{10\pi}{5\pi}=\qty{2}{m/s}\]
（2）位相は $10\pi\times0.05-5\pi\times0.05=0.5\pi-0.25\pi=0.25\pi$ なので、
\[y=0.02\sin\frac{\pi}{4}=0.02\times0.707\approx\qty{0.014}{m}\]
（3）\[A\omega=0.02\times10\pi\approx\qty{0.63}{m/s}\]
\phtashikame
$c=f\lambda=5\times0.4=\qty{2}{m/s}$ で、（1）と合います。\\
（3）の\qty{0.63}{m/s}は、波の速さ\qty{2}{m/s}とは別の量です。振幅を2倍にすると（3）は2倍になりますが、$c$ は変わりません。
\kotae （1）\qty{0.02}{m}、約\qty{31}{rad/s}、約\qty{16}{rad/m}、\qty{5}{Hz}、\qty{0.4}{m}、\qty{2}{m/s}（2）約\qty{0.014}{m}（3）約\qty{0.63}{m/s}
\end{phreidai}
\begin{jibun}
$y=0.05\sin(4\pi t-2\pi x)$（m・s）の波の振動数・波長・速さと、$x=\qty{0.25}{m}$・$t=\qty{0.5}{s}$ のずれ、媒質の速さの最大値を求めましょう。
\bkjans{$f=\qty{2}{Hz}$、$\lambda=\qty{1}{m}$、$c=\qty{2}{m/s}$。位相 $2\pi-0.5\pi=1.5\pi$ で $y=\qty{-0.05}{m}$。$A\omega=0.05\times4\pi\approx\qty{0.63}{m/s}$}
\end{jibun}

\label{ph:波-5:ex4}
\begin{phreidai}[進むパルスと、1つの場所のゆれ]
図\ref{ph:波-5:fig-pulse}のパルス（$x=0$ から\qty{0.3}{m}まで高さ0から\qty{0.1}{m}へ上がり、\qty{0.4}{m}で0にもどる）が、速さ\qty{2}{m/s}で右へ進みます。（1）$t=\qty{0.5}{s}$ に、山の頂上はどこにありますか。（2）$x=\qty{2.0}{m}$ の点が動き始める時刻・いちばん高くなる時刻・止まる時刻を求め、その点の $y$-$t$ グラフをかきましょう。
\phbamen
形を保って進む波（③）で、ある場所に波が届いてから通りすぎるまでのゆれを調べる場面です。
\phmushi
パルスは形を変えずに進む（分散も弱まりもない）とします。
\phzu
はじめの形 $y_0(x)$：前の端（右）$x=\qty{0.4}{m}$、頂上 $x=\qty{0.3}{m}$、後ろの端（左）$x=0$。速さ $c=\qty{2}{m/s}$。
\phriyu
$y=y_0(x-ct)$ なので、形のどの点も、$t$ 秒で $ct$ だけ右へ移ります。ある点に届く時刻は、（その点までの距離）÷ $c$ です。
\phtateru
\[\text{頂上の位置}=0.3+2t,\qquad\text{届く時刻}=\frac{2.0-(\text{はじめの位置})}{2}\]
\phtoku
（1）$0.3+2\times0.5=\qty{1.3}{m}$。\\
（2）前の端：$(2.0-0.4)\div2=\qty{0.80}{s}$ に動き始める。\\
頂上：$(2.0-0.3)\div2=\qty{0.85}{s}$ にいちばん高くなる。\\
後ろの端：$(2.0-0)\div2=\qty{1.00}{s}$ に止まる。\\
$y$-$t$ グラフは、0.80秒から0.85秒で急に上がり、1.00秒までゆっくり下がる形です。
\par\smallskip\noindent\begin{tikzpicture}
\begin{axis}[bkplot, width=96mm, height=32mm, xmin=0.6, xmax=1.15, ymin=0, ymax=0.13, xtick={0.8,0.85,1.0}, xticklabels={0.80,0.85,1.00}, ytick={0.1}, yticklabel style={/pgf/number format/fixed}, xlabel={$t$（s）}, ylabel={$y$（m）}, axis x line*=bottom]
  \addplot[bkcurve, line width=1.3pt] coordinates {(0.6,0) (0.8,0) (0.85,0.1) (1.0,0) (1.15,0)};
\end{axis}
\end{tikzpicture}
\phtashikame
$y$-$t$ グラフは、$y$-$x$ グラフ（なだらかに上がって、急に下がる）を左右に裏返した形です。先に届くのはパルスの前の端（急な側）だからです。グラフの横の軸が場所か時刻かで、形が裏返ることに注意します。
\kotae （1）$x=\qty{1.3}{m}$（2）0.80秒に動き始め、0.85秒にいちばん高く、1.00秒に止まる
\end{phreidai}
\begin{jibun}
例題4のパルスが、反対に左へ\qty{2}{m/s}で進むとき（$y=y_0(x+ct)$）、$x=\qty{-1.0}{m}$ の点が動き始める時刻・いちばん高くなる時刻・止まる時刻を求めましょう。
\bkjans{後ろの端（$x=0$）が先に届く。動き始め $1.0\div2=\qty{0.50}{s}$、頂上 $1.3\div2=\qty{0.65}{s}$、止まる $1.4\div2=\qty{0.70}{s}$（ゆっくり上がって、急に下がる）}
\end{jibun}

\label{ph:波-5:ex5}
\begin{phreidai}[2つのパルスの重ね合わせ]
長いひもの上を、高さ\qty{2}{cm}・底の幅\qty{4}{cm}の三角形の山Pが右へ、深さ\qty{1}{cm}・底の幅\qty{4}{cm}の三角形の谷Qが左へ、どちらも\qty{1}{cm/s}で進みます。時刻0に、Pの頂上は $x=\qty{2}{cm}$、Qのいちばん深い所は $x=\qty{10}{cm}$ にあります。（1）時刻\qty{4}{s}のひもの形を求めましょう。（2）時刻\qty{8}{s}の形はどうなりますか。
\phbamen
向かい合って進む2つのパルスが重なる場面です（⑥）。
\phmushi
ひもの曲げのこわさや空気の抵抗を無視し、パルスは形を変えずに進むとします。ゆれは小さく、重ね合わせの原理が使えるとします。
\phzu
$x$〔cm〕で、時刻 $t$ のPの頂上は $2+t$、Qの底は $10-t$。Pのずれ $y_{\mathrm P}$、Qのずれ $y_{\mathrm Q}$。
\phriyu
重ね合わせの原理（⑥）で、ひものずれは $y=y_{\mathrm P}+y_{\mathrm Q}$ です。
\phtateru
$t=\qty{4}{s}$ では、どちらも中心が $x=\qty{6}{cm}$。中心からの距離を $d$（$|d|\leqq2$）とすると、
\[y_{\mathrm P}=2-|d|,\qquad y_{\mathrm Q}=-\Bigl(1-\frac{|d|}{2}\Bigr)\]
\phtoku
（1）\[y=y_{\mathrm P}+y_{\mathrm Q}=2-|d|-1+\frac{|d|}{2}=1-\frac{|d|}{2}\]
$x=4$〜\qty{8}{cm}で、まん中（\qty{6}{cm}）の高さ\qty{1}{cm}の三角形の山です。
\par\smallskip\noindent\begin{tikzpicture}[x=1mm, y=1mm]
  \draw[bkInk, line width=0.5pt, -{Stealth[length=1.6mm]}] (0,0) -- (96,0) node[right, font=\footnotesize] {$x$（cm）};
  \foreach \xx in {0,2,...,12} {\draw[bkInk, line width=0.4pt] ({\xx*7},0) -- ({\xx*7},-1.2); \node[font=\footnotesize, anchor=north] at ({\xx*7},-5) {\xx};}
  \draw[bkBlue!50, line width=0.6pt, dash pattern=on 2pt off 1pt] (28,0) -- (42,8) -- (56,0);
  \draw[bkOrange!70, line width=0.6pt, dash pattern=on 2pt off 1pt] (28,0) -- (42,-4) -- (56,0);
  \draw[bkInk, line width=1.3pt] (0,0) -- (28,0) -- (42,4) -- (56,0) -- (90,0);
  \node[font=\sffamily\footnotesize, anchor=west] at (58,6) {$t=\qty{4}{s}$（太い線が和）};
\end{tikzpicture}\par
（2）$t=\qty{8}{s}$ では、Pの頂上は $x=\qty{10}{cm}$、Qの底は $x=\qty{2}{cm}$。2つはもう重なっていないので、Pはもとの高さ\qty{2}{cm}の山、Qはもとの深さ\qty{1}{cm}の谷にもどっています。
\phtashikame
まん中（$d=0$）で $2+(-1)=1$、両はし（$|d|=2$）で $0+0=0$ になり、和の式と合います。通りすぎた後に形がもどるのは、波の独立性（⑥）のとおりです。
\kotae （1）$x=4$〜\qty{8}{cm}の、高さ\qty{1}{cm}の三角形の山（2）高さ\qty{2}{cm}の山（頂上 $x=\qty{10}{cm}$）と、深さ\qty{1}{cm}の谷（底 $x=\qty{2}{cm}$）
\end{phreidai}
\begin{jibun}
例題5のQが、谷ではなく高さ\qty{1}{cm}の山だったら、時刻\qty{4}{s}のひもの形はどうなりますか。
\bkjans{$y=(2-|d|)+(1-|d|/2)=3-1.5|d|$。$x=4$〜\qty{8}{cm}の、高さ\qty{3}{cm}の三角形の山}
\end{jibun}

\label{ph:波-5:ex6}
\begin{phreidai}[初期微動から震源までの距離]
地震のとき、ある地点で、P波が届いてからS波が届くまでに\qty{10}{s}かかりました（初期微動が10秒続いた）。P波の速さを\qty{6.0}{km/s}、S波を\qty{3.5}{km/s}として、震源までの距離を求めましょう。
\phbamen
1つの震源から同時に出た2つの波が、違う速さで同じ地点へ届く場面です（①）。
\phmushi
波はまっすぐ進み、どちらも速さは一定とします（本当は、深さで速さが変わり、道すじが曲がる）。
\phzu
震源までの距離 $D$〔km〕、P波・S波が届くまでの時間 $D/6.0$・$D/3.5$。
\phriyu
時間＝距離÷速さ。2つの時間の差が、初期微動の時間です。
\phtateru
\[\frac{D}{3.5}-\frac{D}{6.0}=10\]
\phtoku
\[D\Bigl(\frac{1}{3.5}-\frac{1}{6.0}\Bigr)=10,\qquad D\times0.1190=10,\qquad D\approx\qty{84}{km}\]
\phtashikame
P波は $84\div6.0=\qty{14}{s}$、S波は $84\div3.5=\qty{24}{s}$ で届き、差は10秒です。$D=8.4\times(\text{初期微動の秒数})$ と、距離は初期微動の時間に比例します（中学理科で学ぶ大森の式）。実際の比例定数は、地域や震源の深さで変わり、6〜8 km/s くらいです。
\kotae 約\qty{84}{km}
\end{phreidai}
\begin{jibun}
初期微動が\qty{7.0}{s}続いたとき、震源までの距離は何 km ですか（P波\qty{6.0}{km/s}、S波\qty{3.5}{km/s}）。
\bkjans{$7.0\div0.1190\approx\qty{59}{km}$}
\end{jibun}

\label{ph:波-5:ex7}
\begin{phreidai}[ギターの弦の張力]
ギターの5弦（ラ、\qty{110}{Hz}）では、弦を伝わる波の速さは\qty{143}{m/s}です（9-1⑨）。この例題は、数学の本と同じ数を見せるので、3けたで書きます。（1）弦の線密度が\qty{6.3}{g/m}のとき、張力を求めましょう。（2）音を少し高くするために、波の速さを1\,\% 大きくするには、張力を何\,\% 大きくすればよいですか。
\phbamen
ペグを回して張力を変え、弦を伝わる波の速さを調節する場面です（④）。
\phmushi
弦の曲げのこわさと重さを無視し、張力は弦のどこでも同じとします。
\phzu
$c=\qty{143}{m/s}$、$\mu=\qty{6.3}{g/m}=\qty{0.0063}{kg/m}$、張力 $F_{\mathrm T}$。
\phriyu
弦の波の速さは $c=\sqrt{F_{\mathrm T}/\mu}$（④）。両辺を2乗すると $F_{\mathrm T}=\mu c^2$ です。
\phtateru
\[F_{\mathrm T}=\mu c^2,\qquad\frac{F_{\mathrm T}'}{F_{\mathrm T}}=\Bigl(\frac{c'}{c}\Bigr)^2\]
\phtoku
（1）\[F_{\mathrm T}=0.0063\times143^2=0.0063\times\num{20449}\approx\qty{129}{N}\]
（2）\[\frac{F_{\mathrm T}'}{F_{\mathrm T}}=1.01^2=1.0201\]
なので、約2\,\%（$129\times0.02\approx\qty{2.6}{N}$）大きくします。
\phtashikame
単位：（kg/m）×（m/s）$^2$ ＝ kg·m/s$^2$ ＝ N。\\
大きさ：\qty{129}{N}は、約\qty{13}{kg}のおもりの重さです。（2）は、$c$ が $\sqrt{F_{\mathrm T}}$ に比例するので、張力の変化の割合は速さの変化の割合の約2倍、と読めます（小さな変化では、2乗が2倍になる。5-17）。
\kotae （1）約\qty{129}{N}（2）約2\,\%
\end{phreidai}
\begin{jibun}
ギターの1弦（\qty{330}{Hz}）の波の速さは約\qty{428}{m/s}です。線密度が\qty{0.57}{g/m}のとき、張力を求めましょう。また、波の速さを2\,\% 大きくするには、張力を何\,\% 大きくすればよいですか。
\bkjans{$0.00057\times428^2\approx\qty{104}{N}$。$1.02^2\approx1.040$ なので、約4\,\%}
\end{jibun}

\label{ph:波-5:ex8}
\begin{phreidai}[レールに耳を当てると]
長い鋼のレールの一方のはしをハンマーでたたきました。\qty{1000}{m}はなれた所で、レールに耳を当てて聞くと、レールを伝わってきた音と、空気を伝わってきた音が、何秒ずれて聞こえますか（鋼は $E=\qty{205}{GPa}$、$\rho=\qty{7850}{kg/m^3}$。空気の音の速さは\qty{343}{m/s}）。
\phbamen
同じ所で出た音が、レール（固体）と空気の2つの道を通って届く場面です（⑤）。
\phmushi
レールを細い棒とみなし、レールの曲がりや、つなぎ目、音の弱まりを無視します。
\phzu
距離 $L=\qty{1000}{m}$、レールの縦波の速さ $c_1=\sqrt{E/\rho}$、空気 $c_2=\qty{343}{m/s}$。
\phriyu
細い棒を伝わる縦波の速さは $\sqrt{E/\rho}$（⑤）。届く時間は距離÷速さです。
\phtateru
\[c_1=\sqrt{\frac{E}{\rho}},\qquad\Delta t=\frac{L}{c_2}-\frac{L}{c_1}\]
\phtoku
\[c_1=\sqrt{\frac{205\times10^9}{7850}}\approx\qty{5110}{m/s}\]
\[\frac{1000}{5110}\approx\qty{0.20}{s},\qquad\frac{1000}{343}\approx\qty{2.92}{s},\qquad\Delta t\approx\qty{2.7}{s}\]
\phtashikame
レールの音は空気の約15倍速く、ほとんどすぐに届きます。昔、遠くの列車の近づく音をレールで聞いたという話は、この速さの違いによるものです（危ないので、まねをしてはいけません）。
\kotae 約\qty{2.7}{s}ずれる（レールが先）
\end{phreidai}
\begin{jibun}
長さ\qty{1}{m}のアルミニウム合金の棒（$E=\qty{69}{GPa}$、$\rho=\qty{2700}{kg/m^3}$）を、縦波がはしからはしまで伝わる時間を求めましょう。
\bkjans{$c=\sqrt{69\times10^9\div2700}\approx\qty{5100}{m/s}$。$1\div5100\approx\qty{2.0e-4}{s}$（\qty{0.20}{ms}）}
\end{jibun}

\label{ph:波-5:ex9}
\begin{phreidai}[水と空気の音の速さ]
（1）水（$K=\qty{2.2}{GPa}$、$\rho=\qty{998}{kg/m^3}$）の中の音の速さを求めましょう。（2）\qty{20}{\celsius}の空気（$p=\qty{101325}{Pa}$、$\rho=\qty{1.20}{kg/m^3}$、$\gamma=1.4$）について、ゆっくり押すとしたとき（$K=p$）と、すばやく押すとしたとき（$K_{\mathrm s}=\gamma p$）の $\sqrt{K/\rho}$ を求め、測った値\qty{343}{m/s}とくらべましょう。
\phbamen
流体の縦波（音）の速さを、体積弾性率と密度から求める場面です（⑤）。
\phmushi
流体のねばりと、音の弱まりを無視します。押し縮められる量は小さいとします。
\phzu
体積弾性率 $K$、密度 $\rho$、速さ $c=\sqrt{K/\rho}$。
\phriyu
流体の縦波の速さは $\sqrt{K/\rho}$（⑤）。気体の $K$ は、押す速さで $p$ と $\gamma p$ に分かれます（力-12・熱-3 ⑧）。
\phtateru
\[c=\sqrt{\frac{K}{\rho}}\]
\phtoku
（1）\[c=\sqrt{\frac{2.2\times10^9}{998}}\approx\qty{1500}{m/s}\]
（2）\[\sqrt{\frac{\num{101325}}{1.20}}\approx\qty{291}{m/s},\qquad\sqrt{\frac{1.4\times\num{101325}}{1.20}}\approx\qty{344}{m/s}\]
\phtashikame
すばやく押す（断熱）とした値が、測った\qty{343}{m/s}とよく合います。ゆっくり押すとした値（約\qty{290}{m/s}）は、ニュートンが1687年に出した値で、約15\,\% 小さすぎました（熱-4 ⑥のコラム）。水の中の音は、空気の約4.4倍速く伝わります。
\kotae （1）約\qty{1500}{m/s}（2）約\qty{290}{m/s}と約\qty{340}{m/s}。断熱の値が測った値と合う
\end{phreidai}
\begin{jibun}
\qty{0}{\celsius}の空気（$\rho=\qty{1.293}{kg/m^3}$、$p=\qty{101325}{Pa}$、$\gamma=1.4$）で、$\sqrt{p/\rho}$ と $\sqrt{\gamma p/\rho}$ を求めましょう。
\bkjans{$\sqrt{\num{101325}\div1.293}\approx\qty{280}{m/s}$、$\sqrt{1.4\times\num{101325}\div1.293}\approx\qty{331}{m/s}$（高校の式 $331.5+0.6\theta$ の $\theta=0$ と合う）}
\end{jibun}

\label{ph:波-5:ex10}
\begin{phreidai}[津波の速さ]
（1）深さ\qty{4000}{m}の海と、深さ\qty{10}{m}の海岸近くで、津波の速さを求めましょう（m/s と km/h で）。（2）深さ\qty{4000}{m}の海を\qty{1000}{km}進むのに、何時間かかりますか。
\phbamen
沖で生まれた津波が、海を進んで海岸へ近づく場面です（⑤）。
\phmushi
津波の波長は深さよりずっと長く、高さは深さよりずっと小さいとします。海の深さは、それぞれの場所で一定とします。
\phzu
$g=\qty{9.8}{m/s^2}$、深さ $h$、速さ $c=\sqrt{gh}$、距離 $L=\qty{1000}{km}$。
\phriyu
波長が深さよりずっと長い水の波なので、浅い水の波の速さ $\sqrt{gh}$（⑤）を使います。
\phtateru
\[c=\sqrt{gh},\qquad t=\frac{L}{c}\]
\phtoku
（1）\[\sqrt{9.8\times4000}\approx\qty{198}{m/s},\quad 198\times3.6\approx\qty{710}{km/h}\]
\[\sqrt{9.8\times10}\approx\qty{9.9}{m/s},\quad 9.9\times3.6\approx\qty{36}{km/h}\]
（2）\[t=\frac{1000\times10^3}{198}\approx\qty{5050}{s}\approx1.4\ \text{時間}\]
\phtashikame
単位：$\sqrt{\text{m/s}^2\times\text{m}}$ ＝ m/s。（m/s）×（\qty{3600}{s/h}）÷（\qty{1000}{m/km}）＝ 3.6 倍で km/h。\\
大きさ：沖ではジェット機ほど、海岸近くでも自転車をこぐ速さより速く、走って逃げ切るのは難しい速さです。
\kotae （1）約\qty{200}{m/s}（約\qty{710}{km/h}）、約\qty{9.9}{m/s}（約\qty{36}{km/h}）（2）約1.4時間
\end{phreidai}
\begin{jibun}
平均の深さ\qty{4000}{m}の太平洋を、津波が\qty{17000}{km}わたるのに、何時間かかりますか。
\bkjans{$\num{17000}\times10^3\div198\approx\qty{8.6e4}{s}$、約24時間（約1日）}
\end{jibun}

\label{ph:波-5:ex11}
\begin{phreidai}[要素に分けた鋼の棒の、いちばん速いゆれ]
衝突のシミュレーション（力-16）で、断面積\qty{1}{cm^2}の鋼の棒を、長さ\qty{5}{mm}の要素に分けます。1つの要素を質量 $m$ のおもりと、ばね定数 $k_{\mathrm s}$ のばねと考えます。（1）$m$・$k_{\mathrm s}$ と、長い波の速さ $c_0$ を求めましょう。（2）いちばん速いゆれの角振動数 $\omega_{\max}$ と、計算がこわれない刻み $\Delta t$ の上限を求めましょう（鋼は $E=\qty{205}{GPa}$、$\rho=\qty{7850}{kg/m^3}$）。
\phbamen
連続した棒を、ばねとおもりの列に置きかえて計算する場面です（⑦）。
\phmushi
棒の横のふくらみ（ポアソン比）を無視し、要素どうしは棒の向きにだけ動くとします。
\phzu
$S=\qty{1e-4}{m^2}$、$a=\qty{0.005}{m}$、$m=\rho Sa$、$k_{\mathrm s}=ES/a$（力-8 ③）。
\phriyu
おもりの列の長い波の速さは $c_0=a\sqrt{k_{\mathrm s}/m}$、いちばん速いゆれは $\omega_{\max}=2\sqrt{k_{\mathrm s}/m}$（⑦）。こわれない条件は $\omega_{\max}\Delta t<2$（力-16 ⑦）。
\phtateru
\[m=\rho Sa,\quad k_{\mathrm s}=\frac{ES}{a},\quad c_0=a\sqrt{\frac{k_{\mathrm s}}{m}}\]
\[\omega_{\max}=2\sqrt{\frac{k_{\mathrm s}}{m}},\quad\Delta t<\frac{2}{\omega_{\max}}\]
\phtoku
（1）\[m=7850\times10^{-4}\times0.005\approx\qty{3.9e-3}{kg}\]
\[k_{\mathrm s}=\frac{205\times10^9\times10^{-4}}{0.005}=\qty{4.1e9}{N/m}\]
\[\sqrt{\frac{k_{\mathrm s}}{m}}=\sqrt{\frac{4.1\times10^9}{3.925\times10^{-3}}}\approx\qty{1.02e6}{rad/s}\]
\[c_0=0.005\times1.02\times10^6\approx\qty{5100}{m/s}\]
（2）\[\omega_{\max}\approx\qty{2.0e6}{rad/s},\qquad\Delta t<\frac{2}{2.04\times10^6}\approx\qty{9.8e-7}{s}\]
\phtashikame
$c_0$ は⑤の $\sqrt{E/\rho}$ と同じです。$\Delta t$ の上限は、波が1つの要素をわたる時間 $0.005\div5100\approx\qty{1e-6}{s}$（力-16 ⑦の約\qty{1}{\micro s}）と同じになりました。$\omega_{\max}$ を振動数にすると約\qty{330}{kHz}で、人には聞こえない、とても速いゆれです。
\kotae （1）$m\approx\qty{3.9}{g}$、$k_{\mathrm s}=\qty{4.1e9}{N/m}$、$c_0\approx\qty{5100}{m/s}$（2）$\omega_{\max}\approx\qty{2.0e6}{rad/s}$、$\Delta t$ は約\qty{1}{\micro s}より短く
\end{phreidai}
\begin{jibun}
アルミニウム合金（$c_0\approx\qty{5100}{m/s}$）の部品を、長さ\qty{1}{mm}の要素に分けるとき、こわれない刻みの上限を求めましょう。
\bkjans{$0.001\div5055\approx\qty{2.0e-7}{s}$（\qty{0.20}{\micro s}）}
\end{jibun}

\label{ph:波-5:ex12}
\begin{phreidai}[ばねとおもりの列の位相速度と群速度]
間かく $a=\qty{0.10}{m}$ で、質量 $m=\qty{0.10}{kg}$ の台車が、ばね定数 $k_{\mathrm s}=\qty{40}{N/m}$ のばねで長くつながった模型があります。（1）長い波の速さ $c_0$ を求めましょう。（2）波長 $4a$ の波の角振動数・位相速度・群速度を求めましょう。（3）波長 $4a$ の波の「かたまり」が\qty{4}{m}進むのに何秒かかりますか。
\phbamen
台車とばねの列を伝わる波の速さが、波長で変わるようすを調べる場面です（⑦）。
\phmushi
台車の車輪の摩擦と、ばねの質量を無視します。ゆれは小さいとします。
\phzu
$\sqrt{k_{\mathrm s}/m}=\sqrt{40/0.10}=\qty{20}{rad/s}$、$\lambda=4a=\qty{0.40}{m}$、$k=2\pi/\lambda$、$ka=\pi/2$。
\phriyu
分散関係 $\omega=2\sqrt{k_{\mathrm s}/m}\,|\sin(ka/2)|$ と、位相速度 $\omega/k$、群速度 $c_0\cos(ka/2)$（⑦）。かたまり（エネルギー）は群速度で進みます。
\phtateru
\[c_0=a\sqrt{\frac{k_{\mathrm s}}{m}},\quad\omega=2\sqrt{\frac{k_{\mathrm s}}{m}}\sin\frac{ka}{2},\quad c_{\mathrm p}=\frac{\omega}{k},\quad c_{\mathrm g}=c_0\cos\frac{ka}{2}\]
\phtoku
（1）$c_0=0.10\times20=\qty{2.0}{m/s}$\\
（2）$k=2\pi\div0.40\approx\qty{15.7}{rad/m}$、$ka/2=\pi/4$ なので、
\[\omega=2\times20\times\sin\frac{\pi}{4}\approx\qty{28.3}{rad/s}\]
\[c_{\mathrm p}=\frac{28.3}{15.7}\approx\qty{1.8}{m/s},\qquad c_{\mathrm g}=2.0\times\cos\frac{\pi}{4}\approx\qty{1.4}{m/s}\]
（3）\[t=\frac{4}{1.41}\approx\qty{2.8}{s}\]
\phtashikame
波長が長いとき（$k\to0$）は、$c_{\mathrm p}$ も $c_{\mathrm g}$ も $c_0=\qty{2.0}{m/s}$ に近づきます（表\ref{ph:波-5:tab-bunsan}）。波長 $4a$ では、どちらも $c_0$ よりおそく、$c_{\mathrm g}<c_{\mathrm p}$ です。かたまりの中の山は、かたまりより速く進み、前のはしで消えて、後ろのはしから現れます。
\kotae （1）\qty{2.0}{m/s}（2）約\qty{28}{rad/s}、約\qty{1.8}{m/s}、約\qty{1.4}{m/s}（3）約\qty{2.8}{s}
\end{phreidai}
\begin{jibun}
$a=\qty{0.05}{m}$、$m=\qty{0.05}{kg}$、$k_{\mathrm s}=\qty{20}{N/m}$ の列で、$c_0$ と、波長 $6a$ の波の位相速度・群速度を求めましょう。
\bkjans{$\sqrt{20/0.05}=20$、$c_0=\qty{1.0}{m/s}$。$ka/2=\pi/6$。$c_{\mathrm p}=c_0\sin(\pi/6)\div(\pi/6)\approx\qty{0.95}{m/s}$、$c_{\mathrm g}=1.0\times\cos(\pi/6)\approx\qty{0.87}{m/s}$}
\end{jibun}

\label{ph:波-5:ex13}
\begin{phreidai}[波の表計算を手で確かめる]
⑧の表計算（はじめの形は点9・10・11 が0.5・1・0.5、手をはなした瞬間は止まっている）で、$s=1.1$ のとき、（1）1歩めの点10と点8、（2）2歩めの点10 の値を、式から手で計算しましょう。（3）$c=\qty{100}{m/s}$、$\Delta x=\qty{0.01}{m}$ の弦で、$s=1$ と $s=1.1$ になる $\Delta t$ を求めましょう。
\phbamen
表計算の数が、本当に1歩の式どおりになっているかを手で確かめ、刻みの大きさを決める場面です（⑧）。
\phmushi
⑧と同じ（弦は点の列。両はしは固定）。
\phzu
$s^2=1.21$。1歩めは⑧の「1歩め」の式、2歩めは1歩の式（$n=1$）。
\phriyu
⑧の1歩の式（8-13⑧）を、1つずつの点に使います。
\phtateru
$s=c\,\Delta t/\Delta x$ を $\Delta t$ について解くと $\Delta t=s\,\Delta x/c$。
\phtoku
（1）点10：\[y^{\,1}_{10}=1+0.605\times(0.5-2\times1+0.5)=1-0.605=0.395\]
点8：\[y^{\,1}_{8}=0+0.605\times(0.5-0+0)\approx0.303\]
（2）点9・11 は $0.5+0.605\times(1-1+0)=0.5$ なので、
\[y^{\,2}_{10}=2\times0.395-1+1.21\times(0.5-2\times0.395+0.5)\]
\[=-0.21+1.21\times0.21\approx0.044\]
（3）\[s=1:\ \Delta t=\frac{1\times0.01}{100}=\qty{1e-4}{s},\qquad s=1.1:\ \Delta t=\qty{1.1e-4}{s}\]
\phtashikame
表\ref{ph:波-5:tab-sheet}の $s=1.1$ の行（1歩めの0.40・0.30、2歩めの0.04）と合います。$\Delta t$ をたった1割大きくしただけで、計算はこわれてしまいます。$\Delta t\leqq\qty{1e-4}{s}$ は、8-13⑧の条件と同じです。
\kotae （1）0.395、約0.303（2）約0.044（3）\qty{1.0e-4}{s}と\qty{1.1e-4}{s}
\end{phreidai}
\begin{jibun}
$s=0.5$ のとき、1歩めの点10 と点9 の値を求めましょう。
\bkjans{$s^2/2=0.125$。点10：$1+0.125\times(-1)=0.875$。点9：$0.5+0.125\times(0+1-1)=0.5$}
\end{jibun}

%% ======================== まちがえやすいところ ========================
\Needspace{11\baselineskip}\bksection{tsumazuki}{まちがえやすいところをチェック}
\begin{tsumazuki}{波といっしょに、物も進む？}
\bkng{海の波が岸へ進むので、うかんだボールも波といっしょに岸へ運ばれる。}
\bkok{媒質は、もとの場所の周りでゆれるだけ。進むのはゆれの形（とエネルギー）。ボールは、その場で上下（と前後）にゆれる。}
ボールが岸へ流れるのは、風や、くだけた波の流れのためです（①）。
\end{tsumazuki}

\begin{tsumazuki}{振動数を上げると、波は速くなる？}
\bkng{スピーカーの振動数を2倍にすると、音は2倍速く伝わる。}
\bkok{波の速さ $c$ は媒質で決まる。振動数を2倍にすると、波長が半分になる（$c=f\lambda$）。}
高い音も低い音も、同じ速さで届くので、遠くの演奏も、音がずれずに聞こえます（②）。
\end{tsumazuki}

\begin{tsumazuki}{媒質の速さと波の速さ}
\bkng{$y=A\sin(\omega t-kx)$ の媒質の点は、波の速さ $c$ で上下に動く。}
\bkok{媒質の点の速さは $A\omega\cos(\omega t-kx)$ で、最大 $A\omega$。波の速さ $c=\omega/k$ とは別の量。}
振幅を大きくすると媒質は速く動きますが、波の速さは変わりません（③・\bkreidairef{3}）。
\end{tsumazuki}

\begin{tsumazuki}{$y$-$x$ グラフと $y$-$t$ グラフ}
\bkng{ある瞬間の波の形（$y$-$x$）と、1つの点のゆれ（$y$-$t$）は、同じ形になる。}
\bkok{横の軸が違う。パルスなら、$y$-$t$ グラフは $y$-$x$ グラフを左右に裏返した形になる。}
先に届くのは、波の前の端だからです（\bkreidairef{4}）。
\end{tsumazuki}

\begin{tsumazuki}{張力を2倍にすると、速さは2倍？}
\bkng{張力を2倍にすると、弦を伝わる波の速さも2倍になる。}
\bkok{$c=\sqrt{F_{\mathrm T}/\mu}$ なので、$\sqrt2\approx1.4$ 倍。速さを2倍にするには、張力を4倍にする。}
平方根の中の量を変えたときは、2乗・平方根に気をつけます（④）。
\end{tsumazuki}

%% ======================== 工学ではこう使う ========================
\Needspace{18\baselineskip}\bksection{phgear}{工学ではこう使う}
\begin{phkougaku}[速さの違う波を使って、早く知らせる・中を調べる]
\textbf{緊急地震速報}：震源の近くの地震計がP波をとらえると、すぐに震源の場所と地震の大きさを計算し、S波（大きなゆれ）が届く前に、テレビやスマホに知らせます。震源から\qty{84}{km}の地点なら、P波は14秒後、S波は24秒後に届くので（例題6）、知らせを受けてから、机の下にもぐるくらいの時間があります。速さの違う2つの波があることを、身を守ることに使っているのです。

\textbf{くいの検査}：建物を支えるコンクリートのくいの頭をハンマーでたたき、縦波がくいの先ではね返って、もどってくるまでの時間をはかります。波の速さがわかっていれば、くいの長さや、途中の割れ目の場所がわかります（はね返りは波-6・波-7）。地面をほらずに、地中のくいを調べられるのです。

\textbf{コンピュータで波を計算する}：車の衝突、建物の地震のゆれ、スピーカーの音など、波が伝わる問題は、物を小さな要素（ばねとおもり）に分けて計算します。そのときは、⑦⑧の刻みの条件（CFL 条件）を守り、要素の大きさは、表したい波の波長の10分の1くらいより小さくします（表\ref{ph:波-5:tab-bunsan}で、$\lambda=10a$ なら速さのずれは約2\,\%）。

\begin{phketa}[波の速さの目安]
\phketaline{空気の音（\qty{20}{\celsius}）}{\qty{343}{m/s}（\qty{1}{km}を約3秒）}
\phketaline{水の音・鋼の縦波}{約\qty{1500}{m/s}・約\qty{5100}{m/s}（棒）〜\qty{5900}{m/s}}
\phketaline{地震のP波・S波（浅い所）}{約\qty{6}{km/s}・約\qty{3.5}{km/s}}
\phketaline{津波（深さ\qty{4000}{m}）}{約\qty{200}{m/s}（時速約700 km）}
\phboxhead{答えが桁で合っているか}
固体＞液体＞気体の順に速いのがふつうです。鋼の中の音が数百 m/s、空気の中の音が数千 m/s のような答えが出たら、$E$ の GPa を Pa に直し忘れたか、平方根を取り忘れています。
\end{phketa}
\end{phkougaku}
\phstation{振動・音響工学,地震工学,非破壊検査,数値シミュレーション}{進行波、波動方程式、波数、位相速度、群速度、分散関係、CFL 条件}

\begin{column}{湖で、水の中の音の速さをはかる}
1826年、コラドンとシュトゥルムは、スイスのジュネーブ湖で水の中の音の速さをはかりました。夜、1そうの船で、水の中に下げた鐘を鳴らすと同時に火薬を光らせ、約\qty{14}{km}はなれた船で、水に入れた筒に耳を当てて待ちます。光が見えてから鐘の音が届くまでの9秒あまりから、1秒に約\qty{1435}{m}と求めました。⑤の $\sqrt{K/\rho}$ の値とも、今の精密な値（\qty{8}{\celsius}の水で約\qty{1439}{m/s}）とも、ほとんど違いません。
\end{column}

%% ======================== 章末問題 ========================
\bkendproblems
\bkproblemgroup{基本}{2}
\bkproblem 空気の中の音の速さを\qty{340}{m/s}とします。振動数\qty{170}{Hz}と\qty{680}{Hz}の音の波長を求めましょう。
\bkpans{$340\div170=\qty{2.0}{m}$、$340\div680=\qty{0.50}{m}$}
\bkproblem 波長\qty{0.80}{m}、周期\qty{0.20}{s}の波の速さと振動数を求めましょう。
\bkpans{$c=0.80\div0.20=\qty{4.0}{m/s}$、$f=1\div0.20=\qty{5.0}{Hz}$}
\bkproblem 雷の光が見えてから\qty{3.0}{s}後に、音が聞こえました。雷までの距離は約何 m ですか（音の速さ\qty{340}{m/s}。光が届く時間は無視する）。
\bkpans{$340\times3.0=\qty{1020}{m}$、約\qty{1.0}{km}}
\bkproblem 初期微動が\qty{15}{s}続きました。P波\qty{6.0}{km/s}、S波\qty{3.5}{km/s}として、震源までの距離を求めましょう。
\bkpans{$15\div0.1190\approx\qty{126}{km}$、約\qty{130}{km}}
\bkproblem 例題5のPとQが、どちらも高さ\qty{1}{cm}・底の幅\qty{4}{cm}の三角形で、Pは山、Qは谷のとき、時刻\qty{4}{s}のひもの形を答えましょう。
\bkpans{$y=(1-|d|/2)-(1-|d|/2)=0$。ひもはまっすぐ（各点は動いている）}
\bkproblem ギターの弦の張力を2倍にすると、弦を伝わる波の速さは何倍になりますか。また、波の速さを半分にするには、線密度を何倍にすればよいですか（張力は同じ）。
\bkpans{$\sqrt2\approx1.4$ 倍。線密度を4倍}
\bkproblem 振動数\qty{880}{Hz}の音の、空気（\qty{343}{m/s}）の中と水（\qty{1500}{m/s}）の中での波長を求めましょう。
\bkpans{空気 $343\div880\approx\qty{0.39}{m}$、水 $1500\div880\approx\qty{1.7}{m}$}
\bkproblem 例題4のパルス（$x=0$〜\qty{0.4}{m}、頂上 $x=\qty{0.3}{m}$）が、速さ\qty{1.0}{m/s}で右へ進みます。$x=\qty{3.0}{m}$ の点が動き始める時刻・いちばん高くなる時刻・止まる時刻を求めましょう。
\bkpans{動き始め $(3.0-0.4)\div1.0=\qty{2.6}{s}$、頂上 $\qty{2.7}{s}$、止まる $\qty{3.0}{s}$}
\bkproblemgroup{チャレンジ}{4}
\bkproblem 深さ\qty{200}{m}の大陸だなを、津波が\qty{50}{km}進みます。津波の速さと、かかる時間（分）を求めましょう。
\bkpans{$\sqrt{9.8\times200}\approx\qty{44}{m/s}$。$\num{50000}\div44.3\approx\qty{1130}{s}$、約19分}
\bkproblem 例題12の台車の列（$c_0=\qty{2.0}{m/s}$、$\sqrt{k_{\mathrm s}/m}=\qty{20}{rad/s}$、$a=\qty{0.10}{m}$）で、波長 $3a$ の波の角振動数・位相速度・群速度を求めましょう。
\bkpans{$ka/2=\pi/3$。$\omega=40\sin(\pi/3)\approx\qty{34.6}{rad/s}$。$k=2\pi\div0.30\approx\qty{20.9}{rad/m}$、$c_{\mathrm p}\approx\qty{1.65}{m/s}$。$c_{\mathrm g}=2.0\cos(\pi/3)=\qty{1.0}{m/s}$}
\bkproblem アルミニウム合金（$E=\qty{69}{GPa}$、$\rho=\qty{2700}{kg/m^3}$）の部品を、長さ\qty{2}{mm}の要素に分けて衝突を計算します。こわれない刻みの上限と、\qty{10}{ms}の衝突を計算するのに要る刻みの回数を求めましょう。
\bkpans{$c=\sqrt{69\times10^9\div2700}\approx\qty{5060}{m/s}$。$\Delta t<0.002\div5060\approx\qty{4.0e-7}{s}$。$0.010\div(4.0\times10^{-7})\approx\num{25000}$ 回以上}
\bkproblem $y=0.03\sin(20\pi t-4\pi x)$（m・s）が、波動方程式 $\partial^2y/\partial t^2=c^2\,\partial^2y/\partial x^2$ をみたすことを確かめ、$c$ を求めましょう。また、振動数・波長と、媒質の速さの最大値を求めましょう。
\bkpans{$\partial^2y/\partial t^2=-(20\pi)^2y$、$\partial^2y/\partial x^2=-(4\pi)^2y$ なので、$c^2=(20\pi)^2\div(4\pi)^2=25$、$c=\qty{5}{m/s}$。$f=\qty{10}{Hz}$、$\lambda=\qty{0.5}{m}$。$A\omega=0.03\times20\pi\approx\qty{1.9}{m/s}$}

\begin{furikaeri}
\begin{checklist}
\item 波で進むのは「ゆれの形」だと、横波・縦波で説明できる。\bkfurref{ph:波-5:sec1}
\item 振幅・波長・周期・振動数と $c=f\lambda$ を使いこなせる。\bkfurref{ph:波-5:sec2}
\item 正弦波の式 $y=A\sin(\omega t-kx)$ からゆれを求められる。\bkfurref{ph:波-5:sec3}
\item 媒質の運動方程式から波動方程式と波の速さを導ける。\bkfurref{ph:波-5:sec4}
\item おもりの列の分散関係・群速度と、計算の刻みの条件がわかる。\bkfurref{ph:波-5:sec7}
\end{checklist}
\end{furikaeri}
\bkfukidashi{sora}{ウェーブも、ギターの弦も、津波も、進んでいるのは「形」なんだね。速さは、もどそうとする力と重さの比で決まるんだ！}
